%% ============================================================
%%  post_submission_report.tex
%%  HypatiaX — Post-Submission Report
%%  Open items stripped from the main paper and supplements during
%%  final pre-submission review. Not part of the JMLR submission.
%%  Started: 2026-09-27
%% ============================================================

\documentclass[11pt]{article}
\usepackage{amsmath,amssymb}
\usepackage{hyperref}
\usepackage{booktabs}
\newcommand{\Rsq}{R^{2}}
\newcommand{\II}{\mathrm{II}}
\title{HypatiaX: Post-Submission Report\\
\large Open items deferred from the JMLR submission}
\author{}
\date{2026-09-27 (living document; last updated 2026-10-08)}

\begin{document}
\maketitle

\noindent This report tracks statements, figures, and pipeline defects that
were identified during final pre-submission review but could not be
confirmed against the released repository (\texttt{sednabcn/LLM-HypatiaX-DEV})
in time for submission. Per policy, none of these were left in the
manuscript as unverified numbers; each was either removed, replaced with an
honest gap statement, or (where the manuscript's number was independently
reproduced) confirmed and left in place — see the closed items log at the
end of this report for those.

\section{Open Items}

\subsection*{Item 1: Missing $K=30$ multi-run LLM instability data (DeFi benchmark)}
\textbf{Status:} Open. \textbf{Risk:} Medium — affects Introduction,
\S10.9 (Stability Under Stochastic Inference), Discussion
Finding 1, one design-principle bullet, the OOD-proxy subsection, and one
case-study paragraph (Michaelis-Menten) in the main paper.

\textbf{What was claimed:} A genuine $K=30$-independent-stochastic-run
instability sweep over the 74-task DeFi benchmark, reporting: median
$\Rsq=1.00$, mean $\Rsq=-0.76$, 6 catastrophic failures under that framing; a
regime-distribution table (two mutually inconsistent versions existed in the
paper's own source materials: 70 tasks/61-2-4-3, and separately 58
tasks/44-2-4-12); a Spearman correlation $\rho=-0.70$ between instability and
extrapolation performance across "70 tasks"; and per-task statistics for
three cases (Portfolio Expected Shortfall for correlated assets: $\mu=0.709$,
$\II=0.311$; Theta of option: $\mu=0.569$, $\II=0.148$, "13 of 30 runs
completed"; Black--Scholes Call Price: $\mu=0.229$, $\II=0.267$, "0 of 28
runs succeeded"); and a Michaelis-Menten case-study instability figure of
$\II=0.31$.

\textbf{What was found:}
\begin{itemize}
  \item No raw multi-run result file for the 74-task DeFi benchmark exists
    anywhere in the repository. The loader's own preferred source,
    \texttt{hypatiax\_defi\_variance\_results.json}, is absent.
  \item Every copy of \texttt{instability\_analysis.csv} in the repository
    (7 copies across \texttt{hypatiax-dev/} and \texttt{hypatiax/} snapshots)
    shows \texttt{n\_runs=1} for every task, not 30.
  \item \texttt{scripts/generate\_tables.py}, function \texttt{gen\_instability()}: when
    the JSON source is missing and it falls back to the CSV, the code sets
    \texttt{"k\_runs": 30} as a hardcoded literal, not read from the data.
    This contradicts the function's own docstring, which claims it "never
    falls back to a hardcoded value." \textbf{This is a pipeline bug, not
    just a data gap}, and will silently mislabel this table again on any
    future rerun until fixed.
  \item The nearest genuinely multi-run data in the repository is 5 seeds
    (42/99/123/777/2024), but only for a separate, much smaller ($\sim$15-task)
    PCA-subset benchmark — not usable as a substitute for the 74-task DeFi
    claim.
  \item A separate file, \texttt{exp1\_instability\_stats.json} (for a
    different, 15-equation Feynman/Chemistry-Biology benchmark, unrelated to
    the DeFi claim), contains its own note: "For full 30-run stochastic
    sweep (Spearman $\rho=-0.70$ claim) see
    \texttt{exp4\_instability\_rf02\_04.ipynb}." That notebook does not
    exist anywhere in the repository.
  \item The repo's own core-15 instability artifact reports
    $\II=0.0$ for Michaelis-Menten, directly contradicting the paper's
    claimed $\II=0.31$ for the same case.
\end{itemize}

\textbf{Action taken in the manuscript:} All of the above figures were
removed from the main paper and replaced with honest gap statements
pointing here. The only instability-adjacent figures retained are those
independently verified from the single-run seed-42 file: median LLM
$\Rsq=1.00$, and 5/74 catastrophic failures ($\Rsq<-10$) at that one setting
(explicitly labelled single-run, not a stability-adjusted estimate).

\textbf{Recommended next steps:}
\begin{enumerate}
  \item \textbf{[DONE]} Fix \texttt{gen\_instability()'s CSV fallback to read \texttt{n\_runs}
    from the data (or skip the table, per the function's own stated policy)
    instead of hardcoding \texttt{k\_runs=30}.} \textbf{Done (2026-09-27)} ---
    confirmed against the live \texttt{scripts/generate\_tables.py} in
    \texttt{github.com/sednabcn/LLM-HypatiaX-DEV} (this was the one
    \texttt{paper\_audit.pdf} pipeline claim that checked out exactly as
    described) and patched; see \texttt{generate\_tables\_PATCHED.py}.
    Not yet merged upstream. This pass did not attempt to locate the
    underlying $K=30$ raw multi-run data itself --- the fix only stops the
    generator from mislabelling a single-run CSV as a 30-run sweep.
  \item If a genuine $K=30$ sweep is re-run, regenerate this entire block
    (Introduction, \S10.9, Discussion Finding 1, the
    design-principle bullet, the OOD-proxy cross-reference, and the
    Michaelis-Menten case study) directly from that raw output.
  \item Locate or re-run \texttt{exp4\_instability\_rf02\_04.ipynb} if the
    Spearman $\rho=-0.70$ result is still wanted for the Feynman/Core-15
    track (separate from the DeFi claim above).
\end{enumerate}

\subsection*{Item 2: Unresolved [VALUE REDACTED] placeholders in the routing
supplement, and a newly-found figure conflict}
\textbf{Status:} Open. \textbf{Risk:} Medium — affects
\texttt{supp\_routing\_improvements.tex} only (Baseline caption,
Projected-Impact-on-the-Leaderboard caption and observation, and the
closing paragraph of that section).

\textbf{What was done (2026-09-27):} A pass was made to resolve all
\texttt{[VALUE REDACTED]} placeholders across the two supplements (11
total: 10 in the routing supplement, 1 in the benchmark report; the main
paper itself contains none — see Item 3 below). 7 of the 11 were
resolvable by direct cross-reference to this paper's own
\S~``The Hybrid Decision-Attribution Bug'' (which already discloses
90.5\,\% [67/74] raw uncorrected, 98.6\,\% [73/74] raw
\texttt{hybrid.success} rate, and 59.5\,\% [44/74] corrected) and have
been filled in with that citation in place.

\textbf{What remains open (2026-09-27 follow-up: itemized per placeholder):}
\begin{itemize}
  \item \textbf{P1 (\texttt{tab:baseline} caption):} ``...disclosed in
    the main paper... as uncorrected for a success-attribution bug.
    [REDACTED] See that section before citing either number as
    current.'' Missing: the bug-corrected $R^2>0.99$ rate
    \emph{specific to this table's $n=66$ legacy population} — not the
    same quantity as the main paper's $59.5\%$, which is corrected on
    $n=74$. \textbf{To close:} re-run the main paper's correction rule
    (\S~``The Hybrid Decision-Attribution Bug'': a success counts only
    if \texttt{hybrid.test\_r2}$>0.99$ \emph{and} the sub-method named
    by \texttt{hybrid.decision} itself succeeded) against whichever raw
    JSON file underlies the $n=66$ ``pre-v3.0 historical'' set, then
    insert the resulting percentage.
  \item \textbf{P2 (\texttt{tab:projected} caption):} ``...gives
    $59.5\%$ (provisional pending a file-provenance ambiguity)... as the
    current figure [REDACTED].'' Missing: confirmation that $59.5\%$
    \emph{is} the current figure, or the corrected value once the
    file-provenance ambiguity is settled. \textbf{To close:} the main
    paper itself flags 3 near-duplicate copies of the seed-42 file with
    slightly different per-task results (58--61\,\% band); pick the
    canonical copy, recompute, and fill this in.
  \item \textbf{P3 (Observation block, after ``sub-method had actually
    failed.''):} Missing: the corrected rate following naturally from
    ``uncorrected $90.5\%$ was inflated by the attribution bug $\to$
    corrected to \_\_\_\%.'' This is structurally the same pattern as
    the already-resolved placeholders around it and \emph{may} simply
    be $59.5\%$ by the same cross-reference logic — it was left pending
    only because it carries the identical ``pending re-verification''
    tag as P1/P2/P4, not because of a structural difference.
    \textbf{To close:} confirm with the original author whether that
    tag here means ``genuinely needs pipeline access'' (same bar as
    P1/P2) or was boilerplate copied from the neighboring sentences; if
    the latter, fill in $59.5\%$ with the same citation used for the
    other 7 resolved placeholders.
  \item \textbf{P4 (closing paragraph, immediately before the
    $74.6\%$/$60.5\%$ conflict):} The messiest of the four.
    \textbf{Working hypothesis (unverified):} $59.5\%$ (main paper) is
    the bug-correction applied to the \emph{raw} v3.0 hybrid with no
    routing fixes; $74.6\%$/$60.5\%$ here may instead be the
    bug-correction applied \emph{after} Fixes~0--5 are layered on (this
    subsection is titled ``Projected Impact on the Leaderboard'' and
    immediately follows the fix-by-fix cumulative table). If so, the
    two figures are not actually contradictory — they describe
    different scenarios and only need a one-line disambiguating label,
    not a number swap. \textbf{To close:} recompute the bug-correction
    against the post-fix result file (if one exists) and check whether
    it lands at $74.6\%/60.5\%$. If it does, add the clarifying label.
    If it does not, one of the two numbers is simply wrong and must be
    corrected against source data.
\end{itemize}

\noindent\textbf{Common blocker for all four:} none is resolvable from
the \LaTeX{} text alone. Each needs
\texttt{hypatiax\_defi\_benchmark\_v3\_results\_seed42.json} (or its
$n{=}66$ legacy counterpart) and ideally \texttt{generate\_tables.py} to
recompute mechanically rather than reason about secondhand; neither was
included in \texttt{files.zip}.

\textbf{2026-09-27 follow-up: blocker partially lifted.} Both files were
retrieved directly from \texttt{github.com/sednabcn/LLM-HypatiaX-DEV}
(\texttt{scripts/generate\_tables.py}, and the canonical seed-42 file at
\texttt{hypatiax-dev/data/results/comparison\_results/noise-noiseless/}
\texttt{noiseless/defi/hypatiax\_defi\_benchmark\_v3\_results\_seed42.json},
md5 \texttt{89037935469176741a5fb4479f2c725f} (as of that date; superseded as canonical by \texttt{cbf503d0\ldots} on 2026-09-29, see Item~5) --- confirmed to match the
hash already cited in Item~C1 below). Note the repository holds
\textbf{three distinct byte-for-byte versions} of this filename across its
data snapshots (md5 \texttt{97f1a16c...} in five other snapshot copies,
\texttt{cbf503d0...} in \texttt{hypatiax/data/...}, and the canonical
\texttt{89037935...} used here) --- independently confirming the
manuscript's own note about near-duplicate seed-42 copies with the
provenance ambiguity flagged in P2.

Mechanically re-running the main paper's correction rule (a success counts
only if \texttt{hybrid.test\_r2}$>0.99$ \emph{and} the sub-method named by
\texttt{hybrid.decision} itself has \texttt{test\_r2}$>0.99$) against this
canonical file gives \textbf{44/74 = 59.5\,\%}, with \textbf{23} uncorroborated
cases broken down as Easy 3/24, Medium 11/29, Hard 9/21 --- an exact,
field-for-field match to both the manuscript's own
\texttt{tab:hybrid-bug-breakdown} and the independent Item~C1
recomputation below. \textbf{This closes P2 and P3}: 59.5\,\% is confirmed
as the current, correct figure by direct mechanical recomputation, not
just by cross-reference, so P3 can be filled in with 59.5\,\% and P2's
file-provenance ambiguity is resolved in favor of the canonical copy above.

\textbf{P1 and P4 remain open.} A search of the retrieved repository did
not turn up an $n=66$ ``pre-v3.0 historical'' result file (the nearest
candidates, \texttt{fixc3\_baseline.json} variants, are a different Gate-C
legacy threshold artifact with $n=90$, not $n=66$ --- though that file
does independently corroborate the already-closed 74/90 legacy-threshold
finding elsewhere in this report), nor a post-fix result file distinguishing
the $74.6\%/60.5\%$ figures in P4 from the $59.5\%$ raw figure. Both
still need either a more targeted repository search (the monorepo has
several parallel data snapshots and legacy/archived directories not
exhaustively enumerated in this pass) or direct input from the original
author on where those two files live.

\textbf{2026-09-27, second pass — partial lead on P4/\texttt{tab:projected},
new open question raised.} All 97 \texttt{\_report.md} files in the
repository were enumerated and checked. One is an exact match for the
per-fix cumulative table's final row: \texttt{hybrid\_pysr/defi/\_report.md}
(present identically in five snapshot copies), experiment \texttt{suppA},
reports \textbf{Hybrid R\textsuperscript{2}$\geq$threshold rate
$=82.2\%$ at $N{=}73$ standard cases (NaN-excluded per method)} ---
matching \texttt{tab:projected}'s ``Fix~4, Cumulative 82.2\%'' row exactly.
\textbf{However, this report's own stated threshold is
$R^2\geq0.80$, not $R^2>0.99$} as \texttt{tab:projected}'s column header
claims for the whole table. This raises a new question this pass cannot
close: either (a) the 82.2\% figure is coincidentally identical under both
thresholds for this run, (b) the table's fix-by-fix column is mislabelled
and is actually reporting the $R^2\geq0.80$ coverage rate rather than
$R^2>0.99$, or (c) a different report at the same path but a different
commit produced 82.2\,\% at the $0.99$ threshold and this pass retrieved
a later/earlier version. \textbf{This should be checked before citing any
row of the per-fix table as an $R^2>0.99$ figure.} No file with
$N{=}66$ (the ``honest'' fixed-denominator figures $72.7\%/69.7\%/75.8\%$
in the same table, or the $66.2\%$ baseline in P1's \texttt{tab:baseline})
was found among any of the 97 reports, nor via targeted search for
\texttt{experiment\_protocol\_defi.py} / \texttt{test\_enhanced\_defi\_}
\texttt{extrapolation.py} output logs. The latter script's own catalogue
comment says ``73 test-case catalogue'' with 6 cases flagged
\texttt{extrapolation\_intractable} (giving 67 standard, not 66) ---
close to but not exactly matching either the $n{=}66$ or $n{=}73$ figures
cited elsewhere, and not resolvable further without an actual run log
(this repository holds the test script but not a stored output from
running it in ``honest,'' NaN-penalized mode). \textbf{Recommendation:}
treat every number in \texttt{tab:projected} as unverified pending (i)
confirmation of which $R^2$ threshold the 82.2\,\% figure actually used,
and (ii) locating or re-running the honest-denominator ($n=66$) pass.

\textbf{2026-09-27, third pass --- full git history searched, dead end
reached, figures removed.} A full clone of the repository (not just the
master-branch snapshot used above) was obtained to search commit history
directly. Finding: \textbf{the entire \texttt{tab:baseline}/}
\texttt{tab:projected} \textbf{leaderboard was introduced in a single
commit} (\texttt{8750ede7}, ``Update paper tex,'' 2026-09-03), as a
brand-new file with no predecessor version and \textbf{no companion data
file committed alongside it}. There is no earlier commit to trace back to.
Separately, the canonical seed-42 JSON was checked for the fields an
``old, in-distribution $R^2$ routing'' baseline would need to be
mechanically replayed; it only stores extrapolation-test results, not the
in-distribution split the old routing rule would have used, so this figure
cannot be reconstructed from data already in hand either.

A repository update was also received the same day (commit
\texttt{901f0fa3}, ``Update data,'' plus a new \texttt{hypatiax\_v4/}
directory with its own \texttt{generate\_tables.py} output). This was
checked in full, including a search restricted specifically to the
\texttt{hypatiax/} directory (as opposed to \texttt{hypatiax-dev},
\texttt{hypatiax-all*}, or \texttt{hypatiax\_v4}) per a direct request to
narrow the search. \textbf{None of it bears on P1 or P4}: it is v4-system
output (already set aside per Item~C4's revert-to-v3 decision, and a
different question --- v4's own headline rate --- from what P1/P4 ask),
and a full-text search of \texttt{hypatiax/} for \texttt{66.2}, \texttt{82.2},
\texttt{74.6}, \texttt{60.5}, \texttt{72.7}, \texttt{69.7} turned up nothing
beyond the already-known \texttt{suppA} report (same threshold mismatch as
before) and one coincidental, unrelated match (an \texttt{error\_pct=82.2\%}
in a Feynman timing log, a different experiment entirely).

\textbf{Decision: figures removed from the manuscript (2026-09-27).}
Given no source file exists for the $n=66$ denominator anywhere in the
repository (current state or history), and the one figure that did match
(82.2\,\%) does so at the wrong $R^2$ threshold, these numbers do not meet
the bar this paper has held everywhere else (a live, checkable source
file). Per the same policy already applied to the CI/SD figures (Item~4
below) and the Five-System Comparison, they have been \textbf{removed
from \texttt{supp\_routing\_improvements\_pached.tex} outright} rather
than left as placeholders: \texttt{tab:baseline}'s Hybrid cell, the entire
\texttt{tab:projected} percentage/status columns (retaining only the
qualitative fix descriptions, which are independently verified via
\S~``Summary of Changes''), the dependent observation/remark blocks, the
Fix-interaction gap-analysis paragraph, and the closing
``HypatiaX leads $74.6\%$ vs.\ $60.5\%$'' claim. \textbf{This closes P1
and P4} --- not by resolving them, but by removing the unverifiable
content they were guarding, which is the correct outcome once a
reasonably exhaustive search (three separate passes: targeted file search,
full \texttt{\_report.md} enumeration, full git history) turns up nothing.
If a source is ever located after all, these can be restored with
attribution at that time.

\textbf{Separately confirmed (2026-09-27):} the ablation-generator concern
raised by \texttt{paper\_audit.pdf} (\S3, Item~3: ``\texttt{gen\_ablation()}
was not updated to match the paper's current nan/-inf-aware format'') does
\emph{not} hold against the live \texttt{scripts/generate\_tables.py}:
its mean-row computation already filters both \texttt{NaN} (via
\texttt{r == r}) and $-\infty$-scale sentinel values (via \texttt{r} $\geq
-10^5$) before averaging. This is a third instance of
\texttt{paper\_audit.pdf} describing a state that does not match the
current codebase (see Item~3 below for the first two). The one open
question this does \emph{not} settle: whether the paper's current table
format wants a Mean row at all (the audit's ``no-aggregate-stats format''
description suggests it may not) --- that is a formatting decision, not a
correctness bug, and is left to the author.

\textbf{Confirmed pipeline bug fixed (2026-09-27):} \texttt{gen\_instability()}'s
CSV-fallback path, as retrieved from the repository, does hardcode
\texttt{"k\_runs": 30} exactly as this report's Item~1 and
\texttt{paper\_audit.pdf} both describe --- this one audit claim checks
out fully against the live code. A patched version reading \texttt{k\_runs}
per-row from the CSV (skipping the table if rows disagree or the field is
absent, per the function's own stated policy) has been prepared;
see \texttt{generate\_tables\_PATCHED.py}. Not yet merged upstream ---
this repository is read-only for this review pass.

\textbf{Recommended next steps:}
\begin{enumerate}
  \item \textbf{[DONE]} Obtain the source JSON files (and \texttt{generate\_tables.py
    if possible) and mechanically recompute P1, P2, and P4 as described
    above.} \textbf{Done for P2/P3} (2026-09-27, see above). P1 and P4 still
    need their specific source files located.
  \item \textbf{[DONE]} Decide P3's status (genuine pipeline dependency vs. boilerplate
    copy) and close it accordingly\textbf{ --- done: genuine, and now closed
    with 59.5\,\% per the mechanical recomputation above.}
  \item Test the P4 fix-layering hypothesis against source data before
    assuming either $59.5\%$ or $74.6\%/60.5\%$ is simply wrong; do not
    cite either pending that check. Still open --- post-fix result file not
    located.
\end{enumerate}

\subsection*{Item 3: The external data-integrity audit
(\texttt{paper\_audit.pdf}) contains figures that do not match the
current manuscript}
\textbf{Status:} Open. \textbf{Risk:} Medium-High — affects confidence in
that audit document as a whole, not just the specific figures below.

\textbf{What was found:} A 2026-09-27 pass cross-checking
\texttt{paper\_audit.pdf} against the actual \LaTeX{} sources found two
material discrepancies:
\begin{itemize}
  \item The audit reports 23 live \texttt{[VALUE REDACTED]} placeholders
    in the main paper (34 total across all three documents). A direct
    search of the main paper as provided finds \textbf{zero} such
    placeholders — only one unrelated descriptive mention of
    ``[VALUES REDACTED]'' in prose, not a table value — and 11 total
    across the two supplements, not 34. The audit's claim that
    Table~\texttt{tab:main\_results}'s HypatiaX row is ``entirely
    blanked'' is also not correct as provided: every cell is populated
    except an intentional \texttt{---} for Mean~$\Rsq$, which is
    explained in the table's own caption as a separate, already-labelled
    open item (cross-referenced to Item~1 of this report), not a
    redaction.
  \item The audit describes the hybrid decision-attribution correction
    as affecting 8 of 74 tasks (Easy 22$\to$21 [1], Medium 20$\to$17 [3],
    Hard 10$\to$6 [4], Overall 52$\to$44 [8]). The manuscript's own
    table (\texttt{tab:hybrid-bug-breakdown}) and Item~C1 of this report
    both independently confirm \textbf{23} uncorroborated cases (Easy
    24$\to$21 [3], Medium 29$\to$16 [11], Hard 21$\to$7 [9], Overall
    74$\to$67$\to$44 [23]) — not 8. This is corroborated twice over (the
    manuscript table and the independent C1 recomputation agree with
    each other and disagree with the audit).
  \item \textbf{2026-09-27 follow-up: the audit's \S3 ``Table-Generation
    Pipeline Coverage'' count also does not match.} The audit reports 56
    total \texttt{tab:...} labels across the three documents. A direct
    enumeration of every \texttt{\textbackslash label\{tab:...\}} in the
    three sources as provided finds \textbf{55}, not 56 (18 main paper +
    10 routing supplement + 27 benchmark report; no duplicates, and no
    tables under an alternate \texttt{table:}/\texttt{tbl:} naming
    convention). Two claims from that section \emph{do} check out
    independently: (i) all three tables the audit names as ``blocked or
    needs a human decision'' — \texttt{tab:baseline},
    \texttt{tab:projected}, \texttt{tab:cost\_accuracy\_tradeoff} — exist
    exactly as named, in the routing supplement, matching Open Item~2's
    scope; and (ii) the claim that none of the 56 (55) tables are
    currently wired via \texttt{\textbackslash input\{\}} is confirmed —
    no \texttt{\textbackslash input} or \texttt{\textbackslash include}
    directive appears anywhere in any of the three documents. The
    audit's finer breakdown within that count (27 generated from real
    JSON data / 21 static / 5 intentionally never automated) could
    \textbf{not} be independently checked, since it depends on
    \texttt{generate\_tables.py}, which was not provided alongside the
    \LaTeX{} sources.
\end{itemize}

\textbf{Recommended next steps:}
\begin{enumerate}
  \item Treat \texttt{paper\_audit.pdf}'s specific figures as unverified
    until reconciled against the manuscript version it was actually run
    on — it may have been produced against an earlier draft.
  \item Re-run or re-issue the audit against the current
    \texttt{\_CLEANED} sources before relying on any of its other counts.
  \item Obtain \texttt{generate\_tables.py} and \texttt{table\_wiring\_plan.md}
    (referenced repeatedly by the audit but not included in
    \texttt{files.zip}) before treating the 27/21/5-way pipeline-coverage
    breakdown in the audit's \S3 as confirmed; only the label count and
    the ``blocked'' table names have been independently verified so far.
\end{enumerate}

\subsection*{Item 4: Pure LLM / Neural MLP Field-Level Reproducibility
Gap (\texttt{tab:main\_results})}
\label{psr:mean-r2}
\textbf{Status:} Resolved 2026-09-28 (see below). \textbf{Risk (was):} Medium — affected the credibility of
two of the three rows in the main paper's headline results table, though
not their reported values directly.

\textbf{Provenance note (2026-09-27):} This item did not exist in this
report until today. The main paper's \texttt{tab:main\_results} caption
contained a reference, \verb|\S\ref{psr:mean-r2}|, to a Post-Submission
Report section that was never actually written — the label
\texttt{psr:mean-r2} did not exist anywhere in this document. That
dangling reference was found during a 2026-09-27 pass checking all
cross-document references in the main paper and both supplements for
consistency (see also Item 3, which found a second such issue: three
uses of \texttt{Appendix~\textbackslash ref\{app:statistical\_tests\}} in
the main paper pointing at a label that exists only in
\texttt{supp\_benchmark\_report.tex} — a different, separately-compiled
document, so \texttt{\textbackslash ref} cannot resolve it regardless of
which label is intended. Both the \texttt{psr:mean-r2} and
\texttt{app:statistical\_tests} references in the main paper have been
converted to plain-text pointers so they display correctly however the
document is compiled; this item exists so the content the first pointer
was gesturing at is actually written down somewhere.)

\textbf{What the main paper says:} The Pure~LLM and Neural~MLP rows in
\texttt{tab:main\_results} ``are as originally published and are
corroborated by other figures cited throughout this paper... they do not
independently reproduce, field-for-field, from the single seed-42 result
file used for the HypatiaX row and the corrections in
\S~``The Hybrid Decision-Attribution Bug.''''

\textbf{What was found (2026-09-28):} The canonical seed-42 file
(md5 \texttt{89037935469176741a5fb4479f2c725f}, the non-canonical file, whose figures (62.2\,\%, $-0.4433$, 6 catastrophic) were later replaced by the canonical-file values 60.8\,\%, $-0.3048$, 5; pulled from
\texttt{github.com/sednabcn/LLM-HypatiaX-DEV}) was obtained and all
three rows recomputed under the paper's own definitions. Result: the
published Pure~LLM and Neural~MLP rows did \textbf{not} reproduce.
\begin{center}
\begin{tabular}{lcc}
\toprule
Field & Previously published & Recomputed from file \\
\midrule
Pure LLM $>0.99$ / $>0.9$ (\%) & 62.2 / 62.2 & 59.5 / 62.2 \\
Pure LLM Mean $R^2$ & $-0.7571$ & $-0.4433$ (58 non-NaN) \\
Neural MLP Median / Mean $R^2$ & $-0.4675$ / $-0.9482$ & $-0.6908$ / $-1.4155$ \\
Neural MLP $>0.99$ / $>0.9$ (\%) & 5.4 / 12.2 & 0.0 / 9.5 \\
Neural MLP Catastrophic & 0 & 2 \\
Pure LLM Median, Catastrophic & 1.0000, 6 & 1.0000, 6 (match) \\
\bottomrule
\end{tabular}
\end{center}
The 62.2\,\% published as Pure~LLM's $>0.99$ rate is its $>0.9$ rate (46/74,
equal to the \texttt{pure\_llm.success} flag rate); its true $>0.99$ rate is
59.5\,\% (44/74), which is the figure Table~\texttt{tab:difficulty} already
used, so the paper contained an internal inconsistency that is now removed.
The earlier claim that the code's mean rule reproduces $-0.7571$ ``to
within $\approx0.02$'' is also withdrawn. The origin of the old Pure~LLM and
Neural~MLP figures is not established (they may come from another seed or
run); they are withdrawn rather than corrected.

\textbf{Action taken:} \texttt{tab:main\_results} rebuilt from the file
(HypatiaX shown both uncorrected and corrected); the NN median ($-0.47
\to -0.69$), Pure~LLM near-perfect rate ($62.2 \to 59.5\,\%$) and NN
near-perfect rate ($5.4 \to 0.0\,\%$) updated in the abstract-adjacent
summary, introduction, related-work, findings and conclusion. \textbf{Status: Resolved.}

\subsection*{Item 5: Section 3 --- table-generator reconciliation (2026-09-28)}
\textbf{Status:} Investigated 2026-09-28; one paper defect fixed (Item~7), the rest
open (code changes and missing source data, none applied to the paper).
\textbf{Risk:} Medium --- would silently change a headline metric if
\texttt{generate\_tables.py} output were \verb|\input{}|-ed.

\textbf{Finding A --- corrected-hybrid fallback bug.}
\texttt{gen\_defi\_main()}'s \texttt{\_corrected\_hybrid\_r2()} returns the
named sub-method's $R^2$ only when it is a finite number, and otherwise
falls back to hybrid's own self-reported $R^2$. On seed 42, Pure~LLM is
\texttt{NaN} on 16 tasks; on 10 of them hybrid claims $R^2>0.99$. The
fallback counts those 10 as successes, giving HypatiaX 54/74 = 73.0\,\%
against the paper's 44/74 = 59.5\,\% (Tables~\texttt{tab:difficulty},
\texttt{tab:hybrid-bug-breakdown}). The paper's rule is the correct one
(an unverifiable claim is not a success). \textbf{Fix:} when the named
baseline is \texttt{NaN}, return \texttt{NaN}, not hybrid's value; only
decisions with no baseline in the schema may fall back to hybrid's own
value. With that change the generator's corrected row is
Median 1.0000, Mean $-0.4637$ (62 non-NaN), $>0.99$ 59.5\,\%,
$>0.9$ 60.8\,\%, Catastrophic 6.

\textbf{Finding B --- file naming.} The current public repository's
generator already uses an un-versioned glob
(\texttt{hypatiax\_defi\_benchmark*results*.json}) that resolves to the
v3 seed-42 file in \texttt{defi/}; no v4 file is selected. Per Item~C4
the paper cites v3 only. Confirm the working copy used for the final
build has no hard-coded \texttt{v4} lookups before wiring any DeFi table.

\textbf{Finding C --- the two earlier spot-checks stand, with one correction.}
\texttt{tab:llm\_ablation} (generator adds a Time column and Mean row, collapses
extreme values, conflates \texttt{nan}/$-\infty$) must not be wired in as-is.
\texttt{tab:nguyen12} is superseded by Item~6 below: the paper's own 12/12
table was itself wrong, so the earlier ``generator = 5-seed mean, paper =
seed-42'' comparison is no longer the right question.

\textbf{Finding D --- classification of the 56 tables (measured, 2026-09-28).}
The three documents carry 55 \verb|\label|\texttt{\{tab:\dots\}} tables; the five-seed Nguyen table is
unlabelled, making 56. None is wired with \verb|\input{}|; every table is hand-placed in the
\texttt{.tex}. \texttt{generate\_tables.py} on \texttt{DEV} master (identical to the earlier
upload) was run against the cloned repository's \texttt{hypatiax/data/results}, and again against
\texttt{hypatiax/data\_old/results}, and each output was compared with the paper's table.
Earlier working counts (27/21/5/3) used category definitions that are not recorded; the split
below is the one measured here.
\begin{center}\small
\begin{tabular}{rl}
\toprule
22 & generator runs and writes a file for the label (detail below) \\
25 & generator has the label but skips it: source missing at the path it reads \\
 8 & no generator (manual, withdrawn or estimated; see below) \\
 1 & five-seed Nguyen table: unlabelled, no generator \\
\bottomrule
\end{tabular}
\end{center}
\emph{The 22 that run:} 6 five-system tables (exp1\_five and exp2\_five, full, performance,
extrapolation) reproduce the paper's medians, means and train $\Rsq$ exactly from
\texttt{hypatiax/data\_old} (not from the current \texttt{hypatiax/data}; see Item~7);
1 (\texttt{tab:nguyen12}) was rebuilt by hand and the generator's own output differs (Item~6);
1 (\texttt{tab:llm\_ablation}) must not be wired in (Finding~C); 2 are label collisions
(\texttt{tab:hyperparameters}: generator emits NN seeds, timeouts and thresholds, the paper
prints PySR population/generation settings; \texttt{tab:noise\_sensitivity}: generator emits an
M3/M4 rate at one $\sigma$, the paper prints a 7-level $\sigma/\mathrm{std}(y)$ table with 75
trials per level, whose source is unidentified); 5 sweep tables (\texttt{r2\_noise},
\texttt{rr\_noise}, \texttt{time\_noise}, \texttt{sc\_metrics}, \texttt{winrate}) run on
incomplete sources (one $\sigma$ level and 30 tests, against the paper's 150/180 pairs) and so
cannot be checked; 7 descriptive tables (figlist, methods, provenance, runtime\_reproducibility,
software\_env, supplementary\_files, sweeps) were not compared cell by cell.
\emph{The 25 skipped:} 8 DeFi tables (\texttt{main\_results}, \texttt{difficulty},
\texttt{hybrid-bug-breakdown}, \texttt{timing\_full}, \texttt{timing\_llm\_routed\_full},
\texttt{llm\_detailed}, \texttt{defi\_detailed}, \texttt{portfolio\_seed\_sweep}), whose seed-42 and
seed-sweep files exist under \texttt{hypatiax\_v4/noise-noiseless-v4/} but not at
\texttt{results/defi}, where the generator looks; 8 noiseless-protocol tables (\texttt{arch},
\texttt{hardcoded}, \texttt{interpolation\_stats}, \texttt{domain\_success\_detailed},
\texttt{nrmse}, \texttt{overall}, \texttt{sota}, \texttt{wilcoxon}); 2 split tables
(\texttt{randomsplit}, \texttt{pcasplit}); 2 suppA tables (\texttt{timing\_breakdown},
\texttt{scalability}); 5 that need editorial or config files that do not exist
(\texttt{jmlr}, \texttt{fix5\_cases}, \texttt{formulas\_detailed},
\texttt{hyperparameters\_complete}, \texttt{validation\_stats}).
\emph{The 8 with no generator:} \texttt{feynman30-legacy} (withdrawn stub),
\texttt{additional\_stats}, \texttt{changes}, \texttt{equation\_prevalence},
\texttt{conceptual\_complexity}, and the routing supplement's \texttt{baseline},
\texttt{projected} and \texttt{cost\_accuracy\_tradeoff}. The generator itself records the first
five as manual by design; the last three are estimates or projections without a result source.

\textbf{Finding E --- consequences.} (1) A matching label does not mean a matching table (two
collisions above). (2) A generator that finds partial data still writes a table: the sweep
tables would silently print one noise level. (3) Which result generation a table comes from
matters, and the paper mixes them (Item~7). (4) The DeFi headline tables cannot be regenerated
until the generator's search path is pointed at \texttt{hypatiax\_v4/noise-noiseless-v4/} (or the
files are copied to \texttt{results/defi}); \textbf{no DeFi number in the paper was changed by this
pass}. \textbf{Not fixable in the manuscript:} all of the above are repository and pipeline
changes; nothing was wired in.

\textbf{Recommended order:} (1) apply the Finding~A fix; (2) regenerate
\texttt{defi\_main.tex} and diff it against the now-updated
\texttt{tab:main\_results} --- they should agree except for formatting;
(3) only then wire that table via \verb|\input{}|. No table is wired in by
this update.

\subsection*{Item 6: Nguyen-12 rebuilt from the CI result files (2026-09-28)}
\textbf{Status:} Resolved in the manuscript; one point open.
\textbf{Risk (was):} High --- the paper's headline Nguyen result did not reproduce.

\textbf{Finding.} The repository holds two generations of Nguyen-12 result files.
\emph{Older} (\texttt{hypatiax-dev/data}, unchanged since 2026-07-28): seed~42
gives 11/12 for both systems (Nguyen-12 at $R^2=0.99953$), the other four seeds
12/12, and HypatiaX's stored expression is identical to PySR-only's on all 60
records. \emph{Current} (\texttt{hypatiax/data}, CI of 2026-09-24/25): seed~42
gives 10/12 for both, and the five-seed totals are HypatiaX 51/60 (85.0\,\%),
PySR-only 49/60 (81.7\,\%). The manuscript's 12/12 for both systems, and its
``11/12 independently derived'', matched neither. The uploaded consolidated
report's per-seed table reproduces exactly from the current files; its
per-task table does not (it sums to 52/60 for both systems). Correct
per-task counts ($H$/$P$ out of 5): N-3 5/4, N-4 5/2, N-5 2/1, N-7 5/3, N-9 4/4,
N-12 0/5, all others 5/5 (report had N-5 $H$=3, N-4 $P$=4, N-7 $P$=4).

\textbf{Action taken:} \texttt{tab:nguyen12} and the summary note rebuilt from the
current files; five-seed table and Nguyen-12 gate-failure analysis added; the old
12/12 solve counts marked superseded in the reproducibility note; the
benchmark supplement's Nguyen item corrected.

\textbf{Open:} (i) The STRICT-vs-FALLBACK smoke test (CI run 36116137035). Its two
JSON files were inspected (the output itself is still not in the repository). Findings, now in
the manuscript: Nguyen-5 was \emph{not} fixed by the flag (\texttt{llm\_gate\_failed=False}
in both runs, 1/8 accepted; the runs differ by LLM re-sampling and different PySR
expressions); Nguyen-12 FALLBACK ($R^2=0.9979974$ extrapolation, $0.9999083$ train, 6.3\,s) is
character-for-character PySR-only's expression, so it is not a HypatiaX recovery and misses the
0.9999 bar; the test ran 200 iterations / 4 populations / 90\,s against 1000 / 30 / 1100\,s in the
sweep, so its numbers (for example PySR-only Nguyen-11 at 0.489 against 1.0) cannot be read
against the sweep. Gate outcomes at seed~42 are not reproducible except Nguyen-12 (Nguyen-9 was
0/8 accepted in the sweep and 3/8 in the smoke test; Nguyen-12 is 0/8 in all 7 runs). The claim
``no rerun of exp3/exp3b is needed'' is therefore unsupported and not adopted. To get a
fallback result: full-budget run with \texttt{SPARSE\_SEED=fallback}. Commit
\texttt{\_smoke\_test\_exp3\_gate\_fallback/} if it is to be cited at all.
(ii) Which generation is canonical should be stated in the repository README; the
manuscript uses the current one for Nguyen-12 and the older one for the five-system tables
(Item~7). (iii) Independent derivation of HypatiaX's expressions is withdrawn, not
re-established. A \texttt{symbolic\_equivalence\_report.csv} does exist in
\texttt{hypatiax/data/results/extrapolation/} (an earlier note said it did not), but it covers
seed~42 only, and its checker does not define PySR's \texttt{cube} and \texttt{square}
operators: 22 of its 48 rows are marked FAIL with a ``name 'cube' is not defined'' evaluation
error although $R^2\ge0.9999$. It must be fixed (add the two functions to the SymPy namespace)
and rerun on all five seeds before it can support anything. (iv) The canonical alias
\texttt{exp3\_nguyen12\_seed\{seed\}.json} (PR~\#59) is overwritten by any later run of the same seed,
including a reduced-budget one; add a guard that refuses to overwrite when
\texttt{niterations}, \texttt{populations} or \texttt{n\_tasks} differ. (v) In the working copy of
\texttt{exp3\_nguyen12\_consolidated.py} the default is now \texttt{SPARSE\_SEED=fallback};
every result file behind the paper records \texttt{sparse\_seed\_mode: off}, so submission
numbers describe ``off''. Any new run must state its mode wherever H is reported.
(vi) \textbf{Fixed 2026-09-29:} \texttt{gen\_nguyen12()} previously averaged $R^2$ across
seeds and let $-\infty$ through its finite check (7/12 P, 9/12 H). It now reads the per-seed raw files in
\texttt{hypatiax/data/results/extrapolation/} (run-indexed duplicates skipped), counts recoveries per seed at
$0.9999$ with $-\infty$/NaN as failures, and emits the seed-42 table (10/12, 10/12) and the five-seed table
(HypatiaX 10.2/12, PySR-only 9.8/12); both reproduce the manuscript's tables. The generated seed-42 table has
short equation names rather than the manuscript's typeset formulas, so the manuscript keeps its hand-typed
version until the formula strings are added to the generator.

\subsection*{Item 7: Five-system tables --- table defect fixed, generation mix open (2026-09-28)}
\textbf{Status:} Resolved 2026-09-29 by decision: \texttt{hypatiax/data} is canonical. The six
five-system tables, the supplement's Mann--Whitney block and the text quoting them were regenerated from it.
\textbf{Risk:} Low now; the \texttt{data\_old}-based figures below are historical.

\textbf{Finding.} The paper's \texttt{tab:five\_systems\_full}, \texttt{\_performance},
\texttt{\_extrapolation} and the three \texttt{\_exp2five} tables reproduce exactly from
\texttt{hypatiax/data\_old/results/five\_systems} (for example NN near-regime extrapolation
$R^2$ mean $-3.476$, Hybrid median 3406.6\,\%). The current \texttt{hypatiax/data} re-run gives different
values throughout (NN $-4.846$, Hybrid median 2415.2\,\% on $n=10$, Pure~LLM train $R^2$
$-3.077$ against $-3.468$). The paper therefore mixes generations: Nguyen-12 from the current
one, five-system from the older one. The value 3406.6 is not stored in any \texttt{five\_systems} file of either
generation; it is computed by the generator.
\textbf{Defect fixed:} \texttt{tab:five\_systems\_full} carried $n=11/15/1/13$ and the standard
deviation of the \emph{train} $R^2$ in its Std column, so every 95\,\% CI collapsed to
$[\text{mean},\text{mean}]$. Rebuilt with the generator from the same file: $n=9/12/1/10$, Std and CI of the
extrapolation-error column (Hybrid $[-8577.5,\,50533.5]$, System~2 $[-947.8,\,24036.6]$). Means and
medians unchanged. Because the intervals overlap, the ordering ``System~2 has the lower mean'' is now
stated as not statistically established.
\textbf{Done 2026-09-29 (canonical = \texttt{hypatiax/data}).} New values: Hybrid v50\_2 median 2415.2\,\% /
mean 52730.3\,\% ($n=10$); NN 13283.9\,\% / 37437.0\,\% ($n=12$); System~2 1253.2\,\% / 4539.4\,\% ($n=10$);
Pure LLM $n=2$ (no usable CI). Consequences: System~2 now has the lower \emph{median} as well as mean; the Hybrid mean
is dominated by one equation at $1.634\times10^{27}\,\%$, so the supplement's Cohen-$d$ values ($d=-0.41$; the
``conservative'' $-4.06\times10^{18}$) are uninterpretable and are flagged as such; Mann--Whitney is $n=12/15$,
$U=74.0$, $p=0.225$ (previously $n=11/15$, $U=68.0$, $p=0.234$), no rejection either way. The generator's
``conservative $d$'' equation printed only $\mu_{NN}$ in its numerator while showing the value of
$(\mu_{NN}-\mu_H)/\sigma_{NN}$; fixed. Huge values now print in scientific notation.
\textbf{Still open:} the supplement's medium-regime block reports a Hybrid mean next to the same median as the main table
without saying they are different quantities.

\subsection*{Item 8: Text removed from the Nguyen reproducibility note (2026-09-28)}
\textbf{Status:} Removed from the manuscript; kept here for the record.

The paragraph ``Reproducibility note: LLM model identity\ldots'' was cut back to what still
stands (model identity, dead \texttt{LLM\_MODEL} variable, unseeded sampling, the three action
items). Removed as depending on the withdrawn 12/12 table: (a) a ``recent seed-123 rerun recovers
12/12 against 11/12 originally documented'' and ``12/12 for every other seed (99, 123, 777,
2024)''; (b) a five-seed independent-derivation cross-check (seed~42 11/12, seed~777 11/12, seeds
99/123/2024 none), which relied on copy flags absent from every current file; (c) a ``Verified, this
pass'' block saying seed~42's 12/12~P, 12/12~H and 11/12 independently-derived figures matched
\texttt{tab:nguyen12}, which is no longer true; (d) a within-seed timestamp window (late April to late May
2026) for a seed-123 discrepancy that no longer exists; (e) a correction note describing the Success row as
12/12 for both systems. Also corrected: the tablenote ``differs textually on all 12 records'' (on the two
gate aborts the stored expression is the placeholder \texttt{FAILED}); the ``single finite-but-lower case'' sentence
(there are four, all within $5.3\times10^{-6}$ of PySR-only and all above the 0.9999 bar).

\subsection*{Item 9: Section 3 closure pass, 2026-09-29}
\textbf{Decisions taken:} \texttt{hypatiax/data} is canonical; the DeFi source is
\texttt{results/comparison\_results/noise-noiseless/noiseless/defi} (v3; v4 is out of scope); the Nguyen raw JSONs are in
\texttt{results/extrapolation/}.
\textbf{Generator changes} (\texttt{scripts/generate\_tables.py}, patched copy delivered): (a) DeFi search path
extended to that folder and to the \texttt{15/} and \texttt{15\_pca/} seed folders; (b) Finding~A fix (a named baseline whose
score is NaN now yields NaN, not hybrid's own value); (c) the difficulty tier is case-normalised (files store
\texttt{easy}, generator keys were \texttt{Easy}, so the tier and bug-breakdown tables were being silently skipped);
(d) \texttt{gen\_nguyen12()} rewritten (Item~6(vi)); (e) effect-size rendering fixed (Item~7).
\textbf{Result:} labelled tables that now generate rose from 22 to 29 of 55; 18 remain skipped (arch, domain\_success\_detailed,
fix5\_cases, formulas\_detailed, hardcoded, hyperparameters\_complete, interpolation\_stats, jmlr, nrmse, overall, pcasplit,
portfolio\_seed\_sweep, randomsplit, scalability, sota, timing\_breakdown, validation\_stats, wilcoxon).
\textbf{DeFi seed-42 figures moved to the canonical file} (md5 \texttt{cbf503d0\ldots}; the paper had used the
\texttt{hypatiax-dev/data} file, md5 \texttt{89037935\ldots}). Unchanged: 67/74 raw, 44/74 corrected (59.5\,\%), 23 uncorroborated
(3/11/9), tier tables. Changed: Pure LLM \texttt{NaN} 18/74 (was 16), Pure LLM mean $-0.3048$ over 56 tasks
(was $-0.4433$ over 58), Pure LLM $>0.9$ 60.8\,\% (was 62.2\,\%), catastrophic 5 (was 6; rate 6.8\,\% not 8.1\,\%),
uncorrected HypatiaX mean $0.8781$ (was $0.8946$), corrected mean $-0.3203$ (was $-0.4637$),
\texttt{hybrid.success} 72/74 (was 73/74). Edited in the paper and both supplements.
\textbf{Not resolvable from the repository, left unchanged in the manuscript:}
(i) \texttt{tab:timing\_full} and \texttt{tab:timing\_llm\_routed\_full}: the paper's rows (v3c seed~42: Pure LLM
9.69/8.85\,s, Hybrid 1.99/1.52\,s, $n=74$ per seed, 370 pooled) match no file found. The canonical
\texttt{15/} and \texttt{15\_pca/} seed files hold only \textbf{5 records each} (seed~42's 74-record file is in
\texttt{noiseless/defi}); \texttt{data\_old/15} has 74 per seed but gives 8.35--8.85\,s for Pure LLM. The generator,
run on canonical data, prints a five-record ``v4'' table. These need either the original files or a rerun of the
74-task five-seed sweep.
(ii) \texttt{tab:llm\_detailed} and \texttt{tab:defi\_detailed}: label collisions; the paper tables are named-test
tables (for example ``Constant product'' at 45\,s) that no result file supports, while the generator emits per-task
rows with unnamed cases and $\approx1$--12\,s times. Treat the paper versions as unverified.
(iii) \texttt{tab:main\_results} in the generator omits the ``uncorrected'' row that the paper shows; the corrected row
and shared columns agree.
(iv) The sweep, noise and sample-complexity tables, \texttt{noise\_sensitivity}, the eight noiseless-protocol tables and the
tables in Item~5's remaining list still lack complete source files.
\textbf{Recommended next:} locate or rerun the 74-task DeFi seed sweep (closes (i)); decide whether to delete or
regenerate the two collision tables (ii) \textbf{[done 2026-09-30, Item~10: both regenerated]}; then add the CI check that fails when a paper table differs from its
generated version.

\subsection*{Item 10: Per-case tables, hardcode audit and unsupported claims (2026-09-30)}
\textbf{Status:} Manuscript edited (main paper and both supplements); generator patched; timing, PCA-split and
sweep tables remain open below. \textbf{Risk:} Low for what was edited; Medium for the timing tables until their
source files are located.
\textbf{Scope and limits:} the result JSONs and the repository's \texttt{paper/tables/} were not available for this
pass. Every check below is either arithmetic on the paper's own numbers, a reading of \texttt{generate\_tables.py}, or a run
of the generator on \emph{synthetic} records (layout and regression testing only; no synthetic value entered the paper).
Nothing here was checked against raw result files.

\textbf{Finding A --- the two collision tables (Item~9(ii)) are replaced.} The hand-typed per-case tables in Supplement~B
(\texttt{tab:llm\_detailed}, \texttt{tab:defi\_detailed}) disagree with themselves. Recomputed from the DeFi table's own
rows: mean validation 92.49 (stated 91.3), mean time 128.26\,s (stated 115.2), and 17 rows (not 16) marked Pass in the
Extrap.\ column (stated ``16/16''); no per-test extrapolation $R^2$ appears anywhere, so ``average extrapolation
$R^2=0.9650$'' has no support; ``23/23 pass'' conflicts with the 44/74 (59.5\,\%) corrected rate. In the LLM table five
cases are marked Pass at $R^2<0.99$ although the stated validation threshold is 0.99, the supplement's own S1 description
says 60\,\% over 40 tests against 67.5\,\% (27/40) in the table, and the classical-versus-DeFi split (19/20 vs 8/20) has
no source because the pinned file is DeFi-only. \textbf{Action:} both tables, the Failure-Mode column, the summary rows,
the ``Key Findings'' paragraph and the ``Extrapolation Tests'' paragraph were removed and replaced by
\verb|% Auto-generated by tables/generate_tables.py
% Source: hypatiax/data/results/comparison_results/noise-noiseless/noiseless/defi/hypatiax_defi_benchmark_v3_results_seed42.json
% Date:   2026-10-08 10:53
% DO NOT EDIT MANUALLY — re-run 'make tables' to regenerate


\begingroup\small
\begin{longtable}{L{5.4cm}rr}
\caption{Pure LLM performance by test case (pinned DeFi benchmark, seed 42). \texttt{---}: the result file holds no numeric $R^2$ for that case.}\label{tab:llm_detailed}\\
\toprule
\textbf{Case} & \textbf{$R^2$} & \textbf{Time (s)} \\
\midrule
\endfirsthead
\toprule
\textbf{Case} & \textbf{$R^2$} & \textbf{Time (s)} \\
\midrule
\endhead
\bottomrule
\endfoot
Value at Risk at 95\% & 1.0000 & 8.23 \\
Value at Risk at 99\% & 1.0000 & 6.84 \\
Partial liquidation amount & 1.0000 & 12.01 \\
Constant product formula & 1.0000 & 6.71 \\
Simple Staking APY & 1.0000 & 5.95 \\
Loan-to-Value & $-1.41\times10^{5}$ & 7.61 \\
Spot price from AMM & 1.0000 & 6.79 \\
LP share percentage & $-1.8739$ & 5.94 \\
Long position unrealized PnL & 1.0000 & 4.97 \\
Short position unrealized PnL & 1.0000 & 4.74 \\
Funding rate cost & 1.0000 & 7.26 \\
Validator commission adjusted & 1.0000 & 5.69 \\
Slashing penalty & --- & 14.26 \\
Protocol reserve accumulation & 1.0000 & 7.09 \\
Leveraged position notional & 1.0000 & 6.84 \\
Cross-margin available balance & 1.0000 & 8.64 \\
Realized PnL for long & 1.0000 & 4.88 \\
LP fee earnings & 1.0000 & 8.81 \\
Multi-day Value at Risk & 1.0000 & 7.02 \\
ES scaling for multi-day & 1.0000 & 7.66 \\
Annualised Portfolio tracking error & 1.0000 & 7.26 \\
Incremental VaR & 1.0000 & 8.15 \\
Reserve ratio & 1.0000 & 6.63 \\
Impermanent loss breakeven fee rate & 1.0000 & 11.84 \\
Liquidation Price Long & --- & 7.56 \\
Liquidation Price Short & --- & 6.56 \\
Constant Product Price Impact & 1.0000 & 7.78 \\
Effective Leverage & $-12.4502$ & 14.12 \\
Borrowing Interest & 1.0000 & 6.67 \\
Compounding Staking Returns & 1.0000 & 7.21 \\
Portfolio Sharpe Ratio & 1.0000 & 5.50 \\
APY calculation & 1.0000 & 8.62 \\
Capital efficiency & --- & 8.40 \\
AMM output amount & 1.0000 & 10.10 \\
Price slippage percentage & $-3.2163$ & 6.97 \\
Information ratio & 1.0000 & 5.33 \\
Utilization rate of DeFi & --- & 6.64 \\
Borrow APY from utilization & 1.0000 & 6.45 \\
Supply APY from borrow & 1.0000 & 6.38 \\
Health factor & 1.0000 & 9.74 \\
Impermanent loss percentage & $-2.1798$ & 7.72 \\
Options Delta & --- & 16.66 \\
Collateral ratio & --- & 7.55 \\
Funding rate cost (extended) & 1.0000 & 7.88 \\
Staking reward for fixed lock-up & $-0.9246$ & 7.13 \\
Expected Shortfall at 95\% & 1.0000 & 7.17 \\
Expected Shortfall at 99\% & 1.0000 & 7.20 \\
Optimal LP Position (Kelly) & --- & 9.26 \\
Concentrated liquidity position width & $-2.59\times10^{8}$ & 11.11 \\
Portfolio VaR for two correlated & 1.0000 & 9.19 \\
Position margin ratio & 1.0000 & 6.19 \\
Concentrated liquidity position width (v2) & 1.0000 & 7.66 \\
Spot price from AMM reserve & 0.9810 & 10.03 \\
Black-Scholes Call Price & --- & 21.86 \\
Black-Scholes Put Price & --- & 16.19 \\
Component ES & $-0.2749$ & 9.61 \\
Gamma of option & --- & 12.63 \\
Vega of option & --- & 9.73 \\
Multi-Collateral LTV & 1.0000 & 8.55 \\
Correlated Portfolio VaR & 1.0000 & 8.41 \\
Impermanent loss in constant product & 1.0000 & 9.08 \\
Constant product formula (multivariate) & 1.0000 & 5.91 \\
Convexity Adjustment & 0.9998 & 18.49 \\
Liquidation price for leveraged long & --- & 10.06 \\
Liquidation price for leveraged short & --- & 6.68 \\
Maximum safe leverage & --- & 7.92 \\
Required collateral & --- & 8.04 \\
Forward price for derivative & 1.0000 & 6.74 \\
Portfolio Expected Shortfall for correlated & --- & 4.99 \\
Call option intrinsic & $-3.5798$ & 4.68 \\
Put-call parity & 1.0000 & 7.33 \\
Simple options moneyness & $-50.5674$ & 4.92 \\
Uniswap V3 virtual & $-9.41\times10^{6}$ & 14.65 \\
Theta of option & --- & 21.12 \\
\end{longtable}
\endgroup
| and \verb|% Auto-generated by tables/generate_tables.py
% Source: hypatiax/data/results/defi/hypatiax_defi_benchmark_v3_results_seed42.json
% Date:   2026-10-02 12:07
% DO NOT EDIT MANUALLY — re-run 'make tables' to regenerate


\begingroup\small
\begin{longtable}{L{5.0cm}rrl}
\caption{HypatiaX performance by test case (pinned DeFi benchmark, seed 42). $R^2$ is that of the sub-method named by the routing decision, not hybrid's self-reported value; \texttt{---}: that sub-method returned no numeric $R^2$. Decisions without an independent baseline (\texttt{ensemble}, \texttt{v4\_residual\_nn}) show hybrid's own value.}\label{tab:defi_detailed}\\
\toprule
\textbf{Case} & \textbf{$R^2$} & \textbf{Time (s)} & \textbf{Decision} \\
\midrule
\endfirsthead
\toprule
\textbf{Case} & \textbf{$R^2$} & \textbf{Time (s)} & \textbf{Decision} \\
\midrule
\endhead
\bottomrule
\endfoot
Value at Risk at 95\% & 1.0000 & 1.34 & \texttt{llm} \\
Value at Risk at 99\% & 1.0000 & 1.08 & \texttt{llm} \\
Partial liquidation amount & 1.0000 & 1.13 & \texttt{llm} \\
Constant product formula & 1.0000 & 1.30 & \texttt{llm} \\
Simple Staking APY & 1.0000 & 1.30 & \texttt{llm} \\
Loan-to-Value & $-1.41\times10^{5}$ & 1.27 & \texttt{llm} \\
Spot price from AMM & 1.0000 & 1.07 & \texttt{llm} \\
LP share percentage & $-1.8739$ & 1.11 & \texttt{llm} \\
Long position unrealized PnL & 1.0000 & 1.26 & \texttt{llm} \\
Short position unrealized PnL & 1.0000 & 1.78 & \texttt{llm} \\
Funding rate cost & 1.0000 & 1.12 & \texttt{llm} \\
Validator commission adjusted & 1.0000 & 1.22 & \texttt{llm} \\
Slashing penalty & --- & 1.46 & \texttt{llm} \\
Protocol reserve accumulation & 1.0000 & 1.29 & \texttt{llm} \\
Leveraged position notional & 1.0000 & 1.05 & \texttt{llm} \\
Cross-margin available balance & 1.0000 & 1.20 & \texttt{llm} \\
Realized PnL for long & 1.0000 & 1.41 & \texttt{llm} \\
LP fee earnings & 1.0000 & 1.30 & \texttt{llm} \\
Multi-day Value at Risk & 1.0000 & 1.33 & \texttt{llm} \\
ES scaling for multi-day & 1.0000 & 1.25 & \texttt{llm} \\
Annualised Portfolio tracking error & 1.0000 & 1.21 & \texttt{llm} \\
Incremental VaR & 1.0000 & 2.09 & \texttt{llm} \\
Reserve ratio & 1.0000 & 1.35 & \texttt{llm} \\
Impermanent loss breakeven fee rate & 1.0000 & 1.59 & \texttt{llm} \\
Liquidation Price Long & --- & 1.18 & \texttt{llm} \\
Liquidation Price Short & --- & 1.28 & \texttt{llm} \\
Constant Product Price Impact & 1.0000 & 1.78 & \texttt{llm} \\
Effective Leverage & $-12.4502$ & 1.87 & \texttt{llm} \\
Borrowing Interest & 1.0000 & 1.26 & \texttt{llm} \\
Compounding Staking Returns & 1.0000 & 1.26 & \texttt{llm} \\
Portfolio Sharpe Ratio & 1.0000 & 1.13 & \texttt{llm} \\
APY calculation & 1.0000 & 1.10 & \texttt{llm} \\
Capital efficiency & --- & 1.68 & \texttt{llm} \\
AMM output amount & 1.0000 & 1.16 & \texttt{llm} \\
Price slippage percentage & $-3.2163$ & 1.27 & \texttt{llm} \\
Information ratio & 1.0000 & 1.18 & \texttt{llm} \\
Utilization rate of DeFi & --- & 1.69 & \texttt{llm} \\
Borrow APY from utilization & 1.0000 & 1.27 & \texttt{llm} \\
Supply APY from borrow & 1.0000 & 1.17 & \texttt{llm} \\
Health factor & 1.0000 & 1.57 & \texttt{llm} \\
Impermanent loss percentage & $-2.1798$ & 1.47 & \texttt{llm} \\
Options Delta & --- & 1.76 & \texttt{llm} \\
Collateral ratio & --- & 2.30 & \texttt{llm} \\
Funding rate cost (extended) & 1.0000 & 1.29 & \texttt{llm} \\
Staking reward for fixed lock-up & $-0.9246$ & 1.39 & \texttt{llm} \\
Expected Shortfall at 95\% & 1.0000 & 1.15 & \texttt{llm} \\
Expected Shortfall at 99\% & 1.0000 & 1.38 & \texttt{llm} \\
Optimal LP Position (Kelly) & 0.0000 & 2.59 & \texttt{nn\_fallback} \\
Concentrated liquidity position width & $-2.59\times10^{8}$ & 1.39 & \texttt{llm} \\
Portfolio VaR for two correlated & 1.0000 & 1.66 & \texttt{llm} \\
Position margin ratio & 1.0000 & 1.38 & \texttt{llm} \\
Concentrated liquidity position width (v2) & 1.0000 & 1.20 & \texttt{llm} \\
Spot price from AMM reserve & 0.9810 & 1.33 & \texttt{llm} \\
Black-Scholes Call Price & $-1.8315$ & 4.21 & \texttt{nn} \\
Black-Scholes Put Price & $-0.6802$ & 3.36 & \texttt{nn} \\
Component ES & $-0.2749$ & 1.16 & \texttt{llm} \\
Gamma of option & $-0.1060$ & 4.37 & \texttt{nn} \\
Vega of option & 0.6285 & 4.61 & \texttt{nn} \\
Multi-Collateral LTV & 1.0000 & 1.39 & \texttt{llm} \\
Correlated Portfolio VaR & 1.0000 & 1.44 & \texttt{llm} \\
Impermanent loss in constant product & 1.0000 & 1.26 & \texttt{llm} \\
Constant product formula (multivariate) & 1.0000 & 1.29 & \texttt{llm} \\
Convexity Adjustment & 0.9998 & 1.39 & \texttt{llm} \\
Liquidation price for leveraged long & --- & 1.32 & \texttt{llm} \\
Liquidation price for leveraged short & --- & 1.64 & \texttt{llm} \\
Maximum safe leverage & --- & 1.34 & \texttt{llm} \\
Required collateral & --- & 1.08 & \texttt{llm} \\
Forward price for derivative & 1.0000 & 1.92 & \texttt{llm} \\
Portfolio Expected Shortfall for correlated & --- & 1.34 & \texttt{llm} \\
Call option intrinsic & $-3.5798$ & 1.29 & \texttt{llm} \\
Put-call parity & 1.0000 & 1.36 & \texttt{llm} \\
Simple options moneyness & $-50.5674$ & 1.38 & \texttt{llm} \\
Uniswap V3 virtual & $-9.41\times10^{6}$ & 1.24 & \texttt{llm} \\
Theta of option & $-0.8034$ & 13.51 & \texttt{nn} \\
\end{longtable}
\endgroup
| (generated from the pinned v3 seed-42
file); \verb|\usepackage{longtable}| added to Supplement~B; the two ``95\,\% classical vs 40\,\% DeFi'' mentions in the
routing supplement were removed. Each removal carries a \texttt{[Correction]} comment in the \texttt{.tex}. The S1 entry no
longer states ``60\,\% across 40 tests'' (the CSV itself was not seen).

\textbf{Finding B --- generator defects found while wiring (all fixed in \texttt{scripts/generate\_tables.py}).}
(a) \texttt{gen\_suppb\_defi\_detailed()} printed hybrid's self-reported \texttt{test\_r2}, the Item~5 Finding~A defect, so a
straight swap would have contradicted \texttt{tab:main\_results}. The corrected-R$^2$ rule was hoisted to module level and both
DeFi tables now share it (\texttt{defi\_main.tex} was byte-identical before and after on a 74-record synthetic
regression; seven targeted decision cases pass). (b) A 74-row \texttt{table[ht]}/\texttt{tabular} overflowed the page
(``Float too large'' by 282\,pt); both are now \texttt{longtable}. (c) NaN printed as the string \texttt{nan}; decision names
such as \texttt{nn\_fallback} and \texttt{v4\_llm} were unescaped and would have broken compilation; both fixed, and values
of magnitude $\ge10^4$ print in scientific notation. (d) The case-name field of the pinned file is undocumented (Item~9(ii)
saw ``unnamed cases''): the generator now tries ten candidate keys and otherwise numbers rows ``Case~$N$'' and prints a
warning listing the top-level keys of record~0. \textbf{Run it once and check stderr for that warning.} (e) Earlier in
this pass: the noiseless Avg Runtime column read only \texttt{time\_s}, but the noiseless shards store \texttt{time}
(column printed \texttt{---}); it now accepts both. This patch was written after the local run above and has not yet been exercised on the real shards. \texttt{gen\_suppb\_sota()} matched exact method keys, so the
\texttt{EnhancedHybridSystemDeFi (core)} row was silently dropped; it now accepts a key plus suffix.

\textbf{Finding C --- Spearman $\rho=0.89$, $p<0.001$ withdrawn.} It appeared in the main paper (LLM-baseline paragraph and the
reproductive-behaviour follow-up), Supplement~B (Correlation Analysis, and again in the removed Key Findings) and the routing
supplement (Equation Prevalence). It is not supported: the only prevalence figures are order-of-magnitude estimates for eight
equations (caption: ``estimated''; the generator records this table as an estimate with no named source); the table has no
per-equation success column; recomputing from the paper's own eight rows (using the now-removed, itself unsourced, per-case
values) gives $\rho=0.59$--$0.63$; and with $n=8$ the exact two-sided $p$ at $\rho=0.89$ is about 0.005 ($p<0.001$ needs
$\rho\ge0.976$). The sentence ``we tested this by searching Common Crawl snapshots'' contradicts the table's own caption.
\textbf{Action:} all four statements now describe a hypothesis that was not tested; the routing subsection is retitled
``Coverage Hypothesis (illustrative estimates)''; \texttt{tab:equation\_prevalence} is kept and labelled as estimated.

\textbf{Finding D --- the seven ``$4.1\times$--$7.7\times$'' runtime statements.} Recomputed from the two timing tables' own cells:
hybrid-over-neural mean ratios recompute within rounding for every row, and pooled means equal the average of the seed
means. All-task per-seed slowdown is $4.82$--$7.72\times$; LLM-routed per-seed mean $4.07$--$5.76\times$ and median
$3.77$--$5.46\times$. The union of mean slowdowns is $4.07$--$7.72\times$, so six of the seven statements (which say
``mean'') are correct as written. The seventh (``\ldots whether measured by mean or median'') is not: including the LLM-routed medians the
range is $3.8\times$--$7.7\times$, and the lower bound was changed. This validates the arithmetic only; the tables themselves are
still unverified (Item~9(i)).

\textbf{Checker run on the unpatched generator (2026-09-29, local).} \texttt{DIFF=12 EXCEPTION=1 MATCH=14 NOSOURCE=24
PARTIAL=5 UNVERIFIED=4}. Diagnoses of the five tables in question (inferred from the generator code and log; not confirmed
against data): \texttt{winrate}, \texttt{time\_noise} and \texttt{noise\_sensitivity} are built from a single sigma level
and a single $n$ (\texttt{load\_sweep\_json} returns only the newest file and never merges shards; the loaded file is named
\texttt{...nshards04.json}); \texttt{overall}, \texttt{sota}, \texttt{wilcoxon} and \texttt{nrmse} read the noiseless
records through a fallback to \texttt{feynman-tests/exp2}, because the DeFi noiseless shards were reported missing; \texttt{nrmse}'s generated table
contains values the paper lacks (for example $6.6\times10^{6}$ and 12916.4), consistent with finite sentinel-scale values being averaged in although
its caption says sentinel overflow is excluded (which column each value belongs to was not checked). Note that the checker is
one-directional (generated numbers must appear in the paper), so a truncated generated table can report MATCH:
\texttt{r2\_noise}, \texttt{rr\_noise} and \texttt{sc\_metrics} passing is weak evidence until the sweep sources are complete.

\textbf{Still open.} (1) \texttt{tab:timing\_full} and \texttt{tab:timing\_llm\_routed\_full} (Item~9(i)): the paper text says all ten
seed files were independently re-checked, outside \texttt{generate\_table1.py}, at 74 tasks each, while the canonical files hold 5 records each;
needs the original files or the 74-task five-seed rerun. If the PCA half proves unsupported, the range is unchanged on the
v3c rows alone (all-task $5.38$--$7.72\times$, LLM-routed mean $4.07$--$5.76\times$, union $4.07$--$7.72\times$), but the ``both
benchmark variants, 740 tasks'' wording would have to change. (2) \texttt{tab:pcasplit}: internally consistent ($3+2+12=17$;
$10+3=13/30=43.3\,\%$; $13+2+15=30$), but its caption says the other 13 per-equation $R^2$ ``were not captured'' while the
generator's table (all equations in the source file) contains values the paper lacks (0.9979, 0.995, 0.5): diff the 17 rescored
rows against the source, then consider wiring the full table. \textbf{[Diffed; see Item~11 Finding~E.]} (3) \texttt{tab:formulas\_detailed}: ground-truth constants are hand-typed;
they need extracting from the benchmark code into \texttt{config/ground\_truth\_formulas.json}. (4) Not touched, same provenance
problem: ``Across 40 tests, 6 equations suffered catastrophic extrapolation failure'' in the main paper's LLM-extrapolation
paragraph, and the ``(40 tests)'' label on the Pure-LLM notebook in Supplement~B. (5) \texttt{nrmse} sentinel cutoff needs a
decision (the paper's threshold is not stated). (6) Complete noise-sweep, sample-complexity and DeFi noiseless shards are needed before
the sweep and noiseless-protocol tables can be judged. (7) After regenerating, remove the two UNVERIFIED entries for
\texttt{tab:llm\_detailed} and \texttt{tab:defi\_detailed} from \texttt{paper/paper\_table\_exceptions.json} (the checker then reports
both as WIRED; against the synthetic tables it did). \textbf{[Done; CI reports both as WIRED, see Item~11.]}

\textbf{Verification performed.} All three documents compile with identical error and page counts before and after the edits
(the residual errors come from a stub standing in for the missing \texttt{jmlr2e.sty}); the withdrawn claims are absent from the
rendered PDFs. \textbf{Recommended next:} run \texttt{generate\_tables.py} without flags and read stderr for the case-name warning;
rebuild; run the checker; then the timing source question, which gates everything else in this list.

\subsection*{Item 11: Paper-table CI check live; file-selection defect fixed (2026-10-01)}
\textbf{Status:} The CI check recommended in Item~9 exists and passes; one generator defect found and fixed
(local result $\ne$ CI result); the remaining open work is data, not code. \textbf{Risk:} Low for what was
changed; the 25 UNVERIFIED tables below are paper numbers with no data behind them.

\textbf{Finding A --- the check.} Job \texttt{paper\_check} in \texttt{.github/workflows/ci\_postprocess.yml}
regenerates every table from the committed result files into a scratch directory
(\texttt{generate\_tables.py} and \texttt{generate\_nguyen12\_fiveseed.py}, using \texttt{results\_root} from
\texttt{config/experiments.yml}) and runs \texttt{check\_paper\_tables.py --strict-nosource} against
\texttt{paper/*.tex} and \texttt{paper/paper\_table\_exceptions.json}. DIFF, PARTIAL and NOSOURCE rows fail the job;
EXCEPTION and UNVERIFIED rows are listed and never fail. Result of run 36777213622 (60 labels):
\textbf{MATCH=22, EXCEPTION=11, UNVERIFIED=25, WIRED=2}; no DIFF, PARTIAL or NOSOURCE; regression suite 23 passed;
the generator reports 2 skipped tables (\texttt{timing\_detail} is one). This supersedes the local counts in Item~5
Finding~D and the unpatched-generator run recorded in Item~10. The rows are:
\begin{itemize}\small
\item \textbf{MATCH (22):} difficulty, hybrid-bug-breakdown, main\_results, methods, nguyen12\_byeq,
  nguyen12\_fiveseed, noise\_sensitivity (see Finding~C), nrmse, overall, provenance, r2\_noise, rr\_noise,
  sc\_metrics, sota, sweeps, wilcoxon, and the six five-systems tables (full, performance, extrapolation, each for
  exp1\_five and exp2\_five).
\item \textbf{WIRED (2):} defi\_detailed, llm\_detailed (pulled in with \verb|\input{}|, cannot drift; Item~10(7) is done).
\item \textbf{EXCEPTION (11):} figlist, figlist\_complete, figlist\_main, figlist\_supplementary,
  runtime\_reproducibility, software\_env, supplementary\_files (descriptive, change every run); hyperparameters
  (label collision); llm\_ablation; nguyen12 (Item~6); pcasplit.
\item \textbf{UNVERIFIED (25):} additional\_stats, arch, baseline, changes, conceptual\_complexity,
  cost\_accuracy\_tradeoff, domain\_success\_detailed, equation\_prevalence, feynman30-legacy, fix5\_cases,
  formulas\_detailed, hardcoded, hyperparameters\_complete, interpolation\_stats, jmlr, portfolio\_seed\_sweep,
  projected, randomsplit, scalability, time\_noise, timing\_breakdown, timing\_full, timing\_llm\_routed\_full,
  validation\_stats, winrate.
\end{itemize}

\textbf{Finding B --- defect: newest file chosen by modification time.} Six places in
\texttt{scripts/generate\_tables.py} picked the ``newest'' result file with \texttt{os.path.getmtime} or
\texttt{st\_mtime} (\texttt{load\_best}, \texttt{load\_noiseless\_tests}, both branches of \texttt{load\_sweep\_json},
the DeFi-73 extrapolation lookup in \texttt{\_vc\_panel\_defi73}, and the portfolio-variance lookup). On a fresh
\texttt{git} checkout (CI, \texttt{fetch-depth: 1}) every mtime is the checkout time, so the choice differs from a
working tree. Observed: \texttt{tab:noise\_sensitivity} was UNVERIFIED 67\,\% locally (0.5 missing) but MATCH 100\,\%
in CI, and \texttt{tab:time\_noise} was missing 7.8 and 6.6 locally but 7.7 and 6.6 in CI. \textbf{Fix:} a helper
\texttt{\_det\_key(p) = (p.name, str(p))} replaces mtime at all six sites; result shards carry a \texttt{YYYYMMDD\_HHMMSS}
stamp, so name order is chronological order on every machine. \textbf{Verified locally:} regenerating before and
after the change differs in exactly one row (\texttt{noise\_sensitivity}: UNVERIFIED 67\,\% $\to$ MATCH 100\,\%);
the after-report and the CI report differ in no row. \textbf{[CI re-run after the push: to be recorded here.]}
A duplicate-record scan showed that all 30 noiseless tests appear, with differing contents, in many shards
(\texttt{exp2}, \texttt{exp2\_pca\_4060} and its \texttt{\_saved} copy, \texttt{noise-sweep}, six
\texttt{sample-complexity} files), so any reader that merges across those roots is order-dependent; the change made
no table differ except the one above, including \texttt{arch} and \texttt{hardcoded}, which read through that loader.

\textbf{Finding C --- \texttt{tab:noise\_sensitivity} MATCH is not verification.} Its exception says the generator emits
an M3/M4 rate at one $\sigma$ while the paper prints a 7-level $\sigma/\mathrm{std}(y)$ table whose source is
unidentified. A MATCH means only that the small generated table is contained in the larger paper table. The checker
applies exceptions only to PARTIAL, DIFF and NOSOURCE rows, so a coincidental MATCH overrides an UNVERIFIED
exception. \textbf{Open decision:} keep as is, or make an exception whose reason starts with \texttt{UNVERIFIED}
win over MATCH so the row reads UNVERIFIED on every machine.

\textbf{Finding D --- limits of the checker (for interpreting the table above).} It is one-directional (generated numbers
must appear in the paper), so a truncated generated table can MATCH. The tolerance is $\max(5\cdot10^{-4},
5\cdot10^{-4}|b|)$, so typos below that size are invisible: the \texttt{tab:pcasplit} Ideal gas cell (paper 0.9999999933,
data 0.9999999993) is treated as equal and can only be fixed by hand. The wording of two exceptions was stale:
\texttt{tab:time\_noise} and \texttt{tab:winrate} cited ``the same incomplete-source problem as \texttt{tab:r2\_noise}'',
although \texttt{r2\_noise} now MATCHes; both now cite the actual cause (Item~5 Finding~D: one $\sigma$ level and 30 tests
against the paper's 150/180 pairs).

\textbf{Finding E --- \texttt{tab:pcasplit}.} The exception records why the row is EXCEPTION at 84\,\%: all 17 paper rows
reproduce from the generator's \texttt{pca\_test\_r2}; the paper omits three equations that fail in the data
(Newton's gravitational force 0.397, kinetic energy $-1.913$, Fermi--Dirac 0.278); the generator's ``Pass $k$/30'' row has no
counterpart in the paper; and the 0.5 that Item~10 listed among values the paper lacks is read from equation names
($E=0.5CV^2$, $KE=0.5mv^2$), not from data. This resolves the diff asked for in Item~10 ``Still open''~(2).

\textbf{Still open before submission.}
(1) \texttt{tab:pcasplit}: correct the Ideal gas cell to 0.9999999993, state in the caption that the column is HDS
held-out $\Rsq$ and how the 17 of 30 rows were chosen, then delete the last sentence of its exception.
(2) \texttt{tab:wilcoxon} caption: it says ``normal approximation with tie correction'', but the printed $p$-values match an
exact test when there are no zero differences or ties; state the method actually used.
(3) The 25 UNVERIFIED tables, by cause: no result file in the repository (timing\_breakdown, scalability, randomsplit,
portfolio\_seed\_sweep, timing\_full, timing\_llm\_routed\_full); generator needs a config file that does not exist
(validation\_stats, jmlr, formulas\_detailed, hyperparameters\_complete, fix5\_cases); legacy tables with no source
(interpolation\_stats, domain\_success\_detailed); hardcoded and arch (the noiseless records read carry no
\texttt{is\_hardcoded} or \texttt{decision}/\texttt{strategy} fields); manual or estimated by design (baseline, projected,
cost\_accuracy\_tradeoff, changes, conceptual\_complexity, equation\_prevalence, additional\_stats, feynman30-legacy);
incomplete sweep sources (time\_noise, winrate); noise\_sensitivity (source unidentified). For each: locate or rerun the source,
or remove the table or its numbers from the paper (policy in the preamble).
(4) Record the CI result of the file-selection fix (Finding~B).

\subsection*{Item 12: Minor issues from the manuscript audit (2026-10-03)}
\textbf{Status:} Open, all minor. \textbf{Risk:} Low. Source: \texttt{manuscript\_audit\_2026-10-03.pdf}, which checked
\texttt{paper\_audit.pdf} and the cleanup review against the three \texttt{\_cleaned.tex} files and the raw result files.
Item~2's description of remaining \texttt{[VALUE REDACTED]} placeholders is out of date: the cleaned sources contain none (they
appear as \texttt{[Removed]} or \texttt{[Unverified]} notes), so Item~2's open points now concern data, not placeholders.
\begin{enumerate}
  \item \textbf{\texttt{tab:arch} caption names the wrong source.} It says ``regenerated from the five noise-sweep shards''.
    The noise-sweep shards carry no \texttt{metadata.decision}. The counts (M3: 140 ensemble / 10 NN; M4: 123 LLM / 27
    ensemble; 150 records each) reproduce exactly from \texttt{protocol\_core\_noiseless\_20260812\_170344.json} plus the four
    \texttt{protocol\_core\_noisy\_*} files ($\sigma=0.0005, 0.001, 0.005, 0.01$). Correct the caption. Those files
    used thresholds $0.95$ ($\sigma=0.0005, 0.001$), $0.995$ and $0.99$, which the caption's ``pool mixes one noiseless run with
    four noisy runs'' should state. This also lets the UNVERIFIED \texttt{arch} entry in Item~11 be closed once the files
    are in the CI result path.
  \item \textbf{Item~11 tables with a source now identified.} \texttt{tab:time\_noise} (M3 7.0--7.8\,s, M4 6.3--6.9\,s per
    equation, M4 $1.1$--$1.2\times$ faster) and the \texttt{tab:winrate} noise column ($0/150$, $28/150$, $122/150$) reproduce
    from the five noise-sweep shards; the \texttt{winrate} sample-complexity column ($0/180$, $8/180$, $172/180$) and
    \texttt{tab:sc\_metrics} reproduce from the six sample-complexity shards. They can leave the incomplete-sweep list in
    Item~11 item~(3) once the shards are committed.
  \item \textbf{Two live \texttt{\textbackslash input} commands in Supplement~B.} \texttt{tables/llm\_detailed.tex} and
    \texttt{tables/defi\_detailed.tex} are read at build time. Item~3 says no \texttt{\textbackslash input} exists, and the
    audit's ``0 of 56 wired'' is also wrong. The build fails unless both files exist next to the source; regenerate them
    (\texttt{gen\_suppb\_llm\_detailed}) before compiling.
  \item \textbf{Four unresolved cross-document references in Supplement~B} will print \texttt{??}: \texttt{sec:nguyen12} and
    \texttt{tab:nguyen12} (both defined only in the main paper), \texttt{sec:findings:item10b} (defined nowhere; it is
    quoted inside a note that says so), and \texttt{sec:five\_systems} (main paper only). The file loads no \texttt{xr}.
    Replace each with plain text, or load \texttt{xr-hyper} for the main paper.
  \item \textbf{Two PCA-directed assertions conflict with the record-gap wording.} The abstract, introduction, conclusion and
    the Runtime Analysis section state that the released DeFi result files do not record whether the single-axis or PCA-directed variant
    was used. Two other passages assert PCA ordering for the DeFi benchmark: the Pipeline Corrections item
    ``v3.0 uses the PCA-directed ordering described above'' and the out-of-distribution metric section (``a held-out test split
    constructed by PCA-directed ordering''). Reword both to say ``the protocol specifies PCA-directed ordering''. The
    Protocol-30 PCA mentions (\texttt{benchmark\_results\_pca\_4060.json}) are separately sourced and need no change.
  \item \textbf{Timing headline scope.} The ``$3.8$--$11.2\times$ slower'' figure covers the five \texttt{v3} files under both
    framings (all tasks $4.77$--$11.23\times$; LLM-routed only $3.84$--$9.39\times$), and ``$370$ task runs'' counts only
    those. The five PCA rows ($4.51$--$6.85\times$, another $370$ runs) lie inside the range but are not named in the
    headline sentence. Add ``and the five PCA files'' or state that the range is for the \texttt{v3} files.
    \texttt{paper\_audit.pdf}'s ``$2.5$--$9.3\times$'' is the withdrawn claim, not a second result.
  \item \textbf{Sample size of the noise sweep.} The \texttt{protocol\_core} noisy files record ``200-sample run'', which
    supports the $n=200$ in the noise-table captions. The noise-sweep shards themselves do not store $n$. The runner patch that
    records \texttt{n\_samples} is not yet committed; until it is, say in the caption that $n=200$ comes from the protocol files.
  \item \textbf{\texttt{gen\_ablation()} still always writes a Mean row.} Confirmed in \texttt{generate\_tables.py}; the paper's
    \texttt{tab:llm\_ablation} has no aggregate row. Do not wire this generator until it matches (as in the audit's open item).
\end{enumerate}
\textbf{Not an issue (checked):} the word ``aggressive'' (describes the single-axis v3c split, not PCA); the Protocol-30 PCA
passages; the arithmetic and orphaned-text problems from the cleanup review (fixed in the sources).

\subsection*{Item 13: Unverified statistics and editorial notes removed from the benchmark supplement (2026-10-04)}
\textbf{Status:} Open. \textbf{Risk:} Low --- the removed material was already flagged
\texttt{[Unverified]} in the supplement and nothing else in the manuscript cited it (checked: no remaining
\texttt{\textbackslash ref} or \texttt{\textbackslash Cref} points at the removed labels).
\textbf{Where:} \texttt{supp\_benchmark\_report.tex}. None of the values below has a source file or script in the
repository; do not restore any of them without one.

\textbf{(a) Interpolation Performance Comparison (statistical-test appendix).} Removed together with its table
(\texttt{tab:interpolation\_stats}) and the ``Key Observation'' paragraph. The Kruskal--Wallis result
($H=42.3$, $p<0.001$) had already been withdrawn in the supplement; the per-system sample sizes were unequal.
\begin{center}\small
\begin{tabular}{lccccc}
\toprule
System & $n$ & Mean $\Rsq$ & Median $\Rsq$ & SD & Range \\
\midrule
System 3 (LLM+Fallback) & 30 & 1.000 & 1.000 & 0.000 & [1.00, 1.00] \\
Hybrid v50\_2 & 15 & 0.931 & 1.000 & 0.147 & [0.58, 1.00] \\
Neural Network & 15 & 0.940 & 0.952 & 0.082 & [0.78, 1.00] \\
Pure PySR & 18 & 0.982 & 0.998 & 0.045 & [0.88, 1.00] \\
\bottomrule
\end{tabular}
\end{center}
Known conflicts: System~3 has $n=15$ in the main paper's regenerated Core-15 table
(\texttt{tab:five\_systems\_performance}), not 30; Pure PySR is absent from that table; the Hybrid v50\_2 and Neural Network
means ($0.931$, $0.940$) differ from the regenerated train-$\Rsq$ means ($0.997$, $0.982$). The rows appear to predate the
regeneration. The removed ``Key Observation'' said System~3 reaches $\Rsq=1.000$ on interpolation ($n=30$) but was not
designed for extrapolation, and that Hybrid v50\_2's median exceeds its mean because of outliers.

\textbf{(b) Domain-Specific Success Rates} (\texttt{tab:domain\_success\_detailed}; success defined as $\Rsq>0.95$). Not
regenerated from result files; only the column averages were checked. The multiples of $33\%$ imply about three tasks per domain.
\begin{center}\small
\begin{tabular}{lccccc}
\toprule
Method & Physics & Chemistry & Biology & DeFi & Economics \\
\midrule
Hybrid v50\_2 & 100\% & 100\% & 67\% & 100\% & 67\% \\
Pure PySR & 100\% & 100\% & 100\% & 67\% & 67\% \\
LLM-Guided PySR & 100\% & 67\% & 67\% & 67\% & 67\% \\
Neural Network & 100\% & 67\% & 33\% & 67\% & 33\% \\
Pure LLM & 100\% & 33\% & 0\% & 33\% & 33\% \\
Avg across methods & 100\% & 73\% & 53\% & 67\% & 53\% \\
\bottomrule
\end{tabular}
\end{center}
The paragraph that followed (no equation-prevalence measurement, so the main paper's reproductive-behaviour hypothesis is
not tested by a correlation) was removed with the subsection. If the main paper still states that hypothesis as tested,
it needs the same caveat.

\textbf{(c) Validation Layer Effectiveness} (\texttt{tab:validation\_stats}). Error counts were not regenerated from
result files; only the percentages and cumulative column were checked for internal consistency.
\begin{center}\small
\begin{tabular}{lccc}
\toprule
Layer & Errors & \% & Cum.\ \% \\
\midrule
1. Dimensional Analysis & 18 & 12 & 12 \\
2. Physics Constraints & 27 & 18 & 30 \\
3. Numerical Stability & 35 & 23 & 53 \\
4. Prediction Accuracy & 70 & 47 & 100 \\
Total & 150 & 100 & --- \\
\bottomrule
\end{tabular}
\end{center}
The removed ``Key Finding'' claimed that single-layer validation (prediction accuracy only) would miss $53\%$ of the errors
caught by earlier layers. Any similar claim elsewhere needs a source.

\textbf{(d) Additional Statistical Tests} (\texttt{tab:additional\_stats}), manually entered, not reproducible:
Hybrid v50\_2 vs Pure PySR, Mann--Whitney $U=89$, $p=0.18$; System~3 vs Hybrid v50\_2, Mann--Whitney $U=156$, $p=0.03$;
Classical vs DeFi (LLM), Fisher's exact $p<0.001$; run-to-run variability, Friedman $\chi^2=34.2$, $p<0.001$. The two
Mann--Whitney rows disagree with the sample sizes in (a): with $n=15$, $18$, $30$, $U=89$ gives $p\approx0.10$ and $U=156$
gives $p\approx0.10$ (two-sided, normal approximation, no tie correction). The Fisher and Friedman rows give no counts,
degrees of freedom or data. Removed with it: a power analysis (power $0.65$ for $d>0.8$ at $n=3$ per domain; power $0.95$ for
$d>0.6$ at $n=15$ per method; no calculation located, and the $n=15$ disagrees with (a)) and the recommendation to raise
per-domain sample sizes to $n\ge10$. The only statistical test in the supplement recomputed from source files is the Hybrid
v50\_2 vs Neural Network extrapolation test (\texttt{app:statistical\_tests}).

\textbf{(e) Two editorial notes.} (i)~The Phase~2 (noisy, $\Rsq>0.995$) heading in the reproduction commands carried a
note that no table reports a noisy-protocol, six-method, 30-equation aggregate at the $0.995$ threshold; the original label
\texttt{tab:feynman30} does not exist in the paper or supplements. Decide whether the heading should point to a specific
table (for example a per-equation noisy breakdown) or stay unreferenced. (ii)~The ``Near-Perfect Recoveries'' paragraph had
a reference to \texttt{sec:results}, which is the main paper's entire Results section, not a listing; it was removed pending
confirmation of the intended target.

\textbf{Needed to close:} the raw per-equation values and scripts behind (a)--(d), or a decision to leave them out
permanently; the two targets in (e).

\subsection*{Item 14: Revision-history notes removed from the main paper (2026-10-04)}
\textbf{Status:} Closed (text moved); one sub-item open. \textbf{Where:} \texttt{jmlr\_paper\_main.tex}. The manuscript
now states the current figure, its source and its caveat in place; the narration of what earlier drafts said was removed.
Nothing below was restored as a result; every figure in the paper is still the one that reproduces from the released files.
\begin{itemize}\small
\item \textbf{Front-matter note} (``Several claims and one full table\ldots withdrawn or corrected during an internal review;
each affected passage is marked in place'') replaced by a one-sentence note pointing here.
\item \textbf{Abstract (Nguyen-12):} removed the statement that the table was rebuilt on 2026-09-28, that it replaces an
earlier 12/12 (100\,\%) figure that does not reproduce, and that 8/12~P, 7/12~H, 4/12 and a never-run Neural MLP baseline
remain withdrawn. The reproduced figures stay: 10/12 each on seed~42; 51/60 vs 49/60 over five seeds.
\item \textbf{Michaelis--Menten (LLM extrapolation):} removed the bracketed correction. The earlier claim (LLM recovers the form on
some runs, $\Rsq=1.000$, collapses on others, $\Rsq<-68$; $\II=0.31$) has no multi-run source, and the one released
instability artifact reports $\II=0.0$ for this equation. The text now says no instability figure is reported for this case.
\item \textbf{Mean-$\Rsq$ definition (metrics list):} the bracketed correction was folded into the definition: NaN entries are dropped,
the denominator is the non-NaN count, not 74, per \texttt{\_generate\_report()} in \texttt{hypatiax\_defi\_benchmark\_v3c.py}.
Withdrawn: Pure~LLM mean $-0.7571$ (re-derived $-0.3048$; previous source file $-0.4433$) and the claim that the two agree to within
$\approx0.02$.
\item \textbf{Five-system section:} removed the ``Issue 4'' note (an earlier Result box had a CI incompatible with the supplement's SD and
was not reproducible). Removed from \texttt{tab:five\_systems\_full} caption: the 2026-09-28 note that the earlier version listed
$n=11/15/1/13$ and used the train-$\Rsq$ SD in the Std column, collapsing every CI to a point.
\item \textbf{\texttt{tab:main\_results} caption:} removed the revision note. Withdrawn: Pure~LLM $62.2\,\%$ for $>0.99$ (that is the $>0.9$
rate; flag rate 46/74 on the previous file, 45/74 = $60.8\,\%$ on the canonical file), Pure~LLM mean $-0.7571$, and Neural~MLP median
$-0.4675$, mean $-0.9482$, $5.4/12.2\,\%$, 0 catastrophic.
\item \textbf{LLM ablation (\texttt{tab:llm\_ablation}):} removed ``rebuilt\ldots replacing the withdrawn version'' and the description of
the withdrawn transcription (finite values on all six \textit{nan}/$-\infty$ cells, including a fabricated $-12.5553$ ``identical collapse''
on Arrhenius and a fabricated $-635$ on Michaelis--Menten; neither is in \texttt{\_merged.json}). The ``both collapse identically''
claim is no longer mentioned in the Arrhenius text. \emph{Superseded (2026-10-08, see Item~16):} the main paper reports a Mann--Whitney comparison on this table under three explicit
conventions for the \textit{nan}/$-\infty$ cells (failed cells ranked worst: $U_H=72$, $p=0.095$; four \textit{nan} equations excluded from H:
$U_H=72$, $p=0.60$; excluded from both: $U_H=48$, $p=0.43$), recomputed from the raw per-equation shard JSON; no mean is reported. The earlier
$U=126$, $p=0.2948$ and ``Run B'' $U=101.5$, $p=0.6833$ are withdrawn ($U=126$ at $n=15$ vs.\ $15$ implies $p\approx0.56$).
\item \textbf{Protocol-30:} removed ``earlier drafts labelled it Feynman-30'' and ``earlier 9/30, 12/30 and 13/30 figures are withdrawn''.
\item \textbf{Nguyen-12:} removed the protocol-note and item~N-3 corrections (the PCA-protocol claim; the Neural MLP ($N$) column and
``MW $P{>}N$'' test, which had no source), the tablenote on the ``11/12 independently derived'' figure and the ``Withdrawn history''
tablenote (8/12~P, 7/12~H; 4/12 caption; 10/12~P / 11/12~H hand-tally; Neural MLP column), the sentence on the earlier
$U=51.0$, $p=0.893$ test (computed on the withdrawn 12/12 table), and the 2026-09-28 note on solve counts (12/12 per seed, 11/12 for seed~123,
a seed-123 rerun and a ``verified'' seed-42 rescoring). The model-identity findings stay. \textbf{Open:} a full-budget rerun recording
\texttt{temperature} and \texttt{n\_candidates} in the result JSON (FIX-N3-ii, item~(ii)); \texttt{run\_all.sh} STEP~7/8 already use
\texttt{\_03.py} with \texttt{--temperature 0.25}.
\item \textbf{Instability section:} removed the ``Per-task instability statistics removed'' paragraph. Its content: a regime-level
qualification; a Spearman correlation $\rho=-0.70$ across ``70 tasks''; an OOD-proxy caveat; per-task multi-run statistics for
Portfolio Expected Shortfall (correlated assets), Theta of option and Black--Scholes Call Price, including ``13 of 30 runs'' and
``zero successful runs in 28 attempts''. All depend on $K=30$ multi-run raw data that is not in the repository. If a multi-run rerun is done,
rebuild the section from its raw output, not from these figures.
\item \textbf{Smaller items:} Arrhenius restored-subsection note; the earlier 40-test LLM-extrapolation claim; the Celeron/Colab\,T4/Kaggle
three-machine comparison (withdrawn run); the Portfolio seed-sweep ``Verified 2026-10-04'' caption (now a plain source line);
the ``renamed from Liquidation price for leveraged long/short'' note; scope notes on the v3c/PCA table labels (the \texttt{v3c} label
and the 54--55 of 74 routing figure are unsupported and should not be cited); the timing appendix's history of an interim six-reconstruction table,
an interim 5-task seed-file report, and a ten-file description.
\item \textbf{``Outstanding Process/Credibility Items (Do Not Cite Until Resolved)''} retitled ``Process and Credibility Notes'' and
shortened. The July~9 CI dashboard is still not citable; the dashboard defect, the 74/90 baseline (not 81/90), the 1.73$\times$ refutation
and the pointer to Supplementary~A \S10 are kept. Moved here: the \texttt{ablation\_paired.json} 21/30 vs 20/30 root cause (one degenerate
Rabi-frequency pair; fixed in \texttt{run\_analysis.py}), the Gate~C \texttt{source\_files} manifest (11 Feynman shard files selected correctly,
one stale pre-filter file trimmed, 74/90 reproduces from the 11), the second unreconciled \texttt{exp2\_run.log}, and the DeFi near-perfect-rate
figures with no locatable source in the repository or its git history.
\end{itemize}

\subsection*{Item 15: Unverified material removed from the routing supplement (2026-10-04)}
\textbf{Status:} Open. \textbf{Where:} \texttt{supp\_routing\_improvements.tex}. None of the values below has a source file; do not restore
any without one.
\begin{itemize}\small
\item \textbf{Case-Level Impact of Fix~5} (\texttt{tab:fix5\_cases}; marked ``Unverified, manually entered''). Removed with its ``net effect:
catastrophic count $3\to1$'' line. Rows (LLM / Hybrid before / Hybrid after): Correlated Portfolio VaR $+1.000$ / $-12.121$ / $\approx+1.000$;
IL in constant product $+1.000$ / NaN / $\approx+1.000$; Capital efficiency NaN / $-3274$ / $\approx+0.827$; Concentrated liquidity $-20.8$ / $-1606$ / $\approx+0.963$.
\item \textbf{Empirical Timing Breakdown} (\texttt{tab:timing\_breakdown}), \textbf{Scalability} (\texttt{tab:scalability}), \textbf{Runtime vs Interpolation Tradeoff}
(\texttt{tab:cost\_accuracy\_tradeoff}) and the \textbf{Parallelisation and Future Speedup} projection. Pure PySR 390.0\,s
(5.2 / 365.8 / 12.3 / 6.7 by stage); LLM-guided hybrid 70.0\,s (8.5 / 2.0 / 3.0 / 52.5 / 4.0); scalability (points: PySR / hybrid, s) 100: 245/48,
500: 390/70, 1000: 521/95, 5000: 1245/215, 10000: 2340/410; tradeoff table: NN 1.7\,s / 0.940, hybrid 70.0\,s / 0.931, PySR 390.0\,s / 0.982 (the
interpolation values duplicate the removed table of Item~13(a)); a projected $\sim8\times$ speedup on 16+ cores and $\sim$1--2\,s per test with GPU fitness
evaluation. No scalability result file exists in \texttt{hypatiax/data/results}. The section is replaced by a short ``Measured Runtime'' paragraph carrying only
figures regenerated from the five 74-task seed files (hybrid $3.8$--$11.2\times$ slower than the NN; NN $0.27$\,s pooled, $0.18$--$0.36$\,s per seed;
68--69 of 74 routed to the LLM).
\item \textbf{Per-fix cumulative percentages and the $n{=}66$ framing} (baseline 66.2\,\% to Fix~4's 82.2\,\%; ``honest $n{=}66$'' rows; projected Fix~5; the
$72.7\,\%$ vs $69.7\,\%$ remark; the closing claim that Hybrid closes or exceeds Pure LLM's $83.9\,\%$ and leads $74.6\,\%$ vs $60.5\,\%$) were already removed on
2026-09-27; this pass removed the bracketed ``[Removed, 2026-09-27]'' narration around them. The only file matching the Fix~4 figure
(\texttt{hybrid\_pysr/defi/\_report.md}, experiment \texttt{suppA}, $N=73$) uses $\Rsq\ge0.80$, not $>0.99$.
\item \textbf{Open items moved out of the ``Remaining Open Issues'' list} (all resolved, not open): the July~9 CI dashboard
(\texttt{notebook\_execution\_status()}, commit \texttt{9d7a833f}, 2026-07-29), the 74/90 baseline-lock threshold fix, and the superseded 1.73$\times$ speedup.
Their current status is in the main paper's ``Process and Credibility Notes''.
\item \textbf{Quick Start commands.} Checked against \texttt{tree\_repro.txt} (2026-10-04). Not found in the repository: \texttt{hypatiax/experiments/comparison/ultimate\_comparative\_suite\_complete\_.py}
(Core~15), \texttt{hypatiax/analysis/statistical\_analysis\_full.py} and \texttt{supplementaries/generate\_figures/generate\_figures.py}; their \texttt{--v2} flags
were never checked. Removed from the Quick Start; also removed: the DeFi command's \texttt{--fixed-denominator 66 --nan-penalty --v2} flags (tied to the withdrawn $n{=}66$
framing; \texttt{test\_enhanced\_defi\_extrapolation.py} exists but its flags were not checked). The Phase~3 command was kept and its path corrected to the file
that exists, \texttt{hypatiax/experiments/benchmarks/run\_comparative\_suite\_benchmark\_v2.py}; its flags (\texttt{--noiseless}, \texttt{--threshold},
\texttt{--nn-seeds}, \texttt{--samples}, \texttt{--method-timeout}, \texttt{--pysr-timeout}) all exist in the script. The script's own docstring and the main paper call it
\texttt{run\_protocol\_benchmark\_core.py}; output files are \texttt{protocol\_core\_*}. Decide whether to rename the file or keep both names documented.
\item \textbf{Not changed, still unverified:} the baseline table (\texttt{tab:baseline}) rows for Pure LLM ($83.9\,\%$ at $>0.99$, 85.7 at $>0.9$, 4 catastrophic) and Neural Net
($1.5\,\%$, 27.9, 6) share the untraceable $n{=}66$ denominator of the removed Hybrid cell; and the reproducibility checklist expected values (Core~15 Mann--Whitney
$U=0$, $p<10^{-6}$; Hybrid DeFi $96.7\,\%$ at $R^2>0.9999$; Cohen's $d=0.95$; 13 figures) were not checked here.
\end{itemize}


\subsection*{Item 16: Data-integrity closure pass (2026-10-08)}
\textbf{Status:} Closed, with the sub-items listed at the end. \textbf{Where:} \texttt{jmlr\_paper\_main.tex},
\texttt{supp\_benchmark\_report.tex}, \texttt{supp\_routing\_improvements.tex}. Every figure below was recomputed from the
released files: the five \texttt{hypatiax\_defi\_benchmark\_v3\_results\_seed<N>.json} files, the four
\texttt{exp1\_ablation\_results\_shard<k>.json} files and \texttt{portfolio\_variance\_seed\_sweep.json}.
\begin{itemize}\small
\item \textbf{Canonical file.} The seed-42 file hashes to \texttt{cbf503d0b3985b4c1bccf7ed18e0316b}; its pure-LLM catastrophic count
($\Rsq<-10$) is 5 (6 is the seed-99 count and the count of the superseded file).
\item \textbf{Cross-seed counts (replaces the 56.8--63.5\,\% / 72.9--75.7\,\% statement).} Pure LLM at $\Rsq>0.99$: 44, 46, 46, 47, 48 of 74
on seeds 42, 99, 123, 777, 2024 ($59.5$--$64.9\,\%$). Raw hybrid: 67/74 ($90.5\,\%$) on every seed. Corrected hybrid equals pure LLM on every
seed; uncorroborated successes 23, 21, 21, 20, 19; corrected gain $0.0$\,pp in every difficulty tier on every seed at both the $>0.99$ and $>0.9$ criteria.
\item \textbf{``0/74 executed'' removed.} Pure LLM returns a finite $\Rsq$ on 59, 60, 61, 60 of 74 tasks on seeds 99, 123, 777, 2024
(57 flagged executed on seed 42). Only the seed-42 records carry an \texttt{executed} field, which probably produced the zeros;
the 20/74 figure for seed 99 does not appear in the released files. The ``multi-seed extension not reportable'' sentences were replaced by the five-seed result.
\item \textbf{Exp-1 provenance.} $-882.9$ does not occur in the seed-sweep JSON or in the four Core-15 ablation shards (15 equations, one run each, no seed
field; the most negative PySR-only far-$\Rsq$ is $-807.7$, Michaelis--Menten; no subset sum of the far values reproduces it). It is withdrawn.
Seed-99 far-$\Rsq$ is $-15.19$ (HypatiaX) and $-1.23$ (PySR-only); the $1.000$ value belongs to the v3c3 DeFi shards (Item~C7).
\item \textbf{Mann--Whitney (Core-15, far $\Rsq$).} From the raw shards: failed cells ranked worst $U_H=72$, $p=0.095$ (exact $0.098$), $r=-0.36$;
four \textit{nan} excluded from H $U_H=72$, $p=0.60$; excluded from both $U_H=48$, $p=0.43$. The earlier $U_H=62$, $p=0.033$ / $62$, $0.28$ / $42$, $0.21$
do not reproduce from the raw files under any natural convention (the nearest, clipping at 0 and rounding to four decimals, gives $62$, $0.031$). PySR-only is
therefore higher under the first convention but not significantly so at the 5\,\% level (one-sided $p=0.048$ for PySR-only higher).
\item \textbf{Text removed from the manuscript (kept here).} Main paper: the account of the three earlier runtime claims ($73\,\%$ reduction; $1.73\times$ / $1.64\times$
speedup; $2.5\times$--$9.3\times$ slower incl.\ PCA seeds with zero LLM-routed tasks), the unsourced $3.0$/$2.7$\,s Neural MLP timings, ``not the 8 cited in earlier drafts'',
the earlier $+3.4$\,pp (Medium) and $-4.8$\,pp (Hard) gains, the earlier Hard-tier gap, the retraction (and re-retraction) of the Portfolio Expected Shortfall task, the paragraph
that left the seed behind the $90.5\,\%$ headline untraced, the Bayesian $P\approx0.76$ note, the $U=126$, $p=0.2948$ note, and the earlier far-$\Rsq=-18.1$ misattribution to seed~42.
Benchmark supplement: the note on an earlier ten-domain list. Routing supplement: the $U=0$, $p<10^{-6}$ expectation.
\item \textbf{Generated tables.} \texttt{tables/llm\_detailed.tex} and \texttt{tables/defi\_detailed.tex} were regenerated from the seed-42 file with
\texttt{generate\_tables.py} (\texttt{gen\_suppb\_llm\_detailed}, \texttt{gen\_suppb\_defi\_detailed}); the benchmark supplement does not build without them.
\item \textbf{Layout and disclosure.} Figures~3 and~4 moved to their first discussion. The Core-15 tier assignment matches \texttt{\_HM\_TIER} in
\texttt{generate\_figures.py} exactly; the paper now describes it only as a fixed row grouping used to order the heatmaps (Item~17, point~1).
\item \textbf{Open sub-items.} (a) The seed-2024 HypatiaX row of the sweep has train $\Rsq=0.9975455190273314$ and the expression
\texttt{sqrt(s1*(2*rho*s2 + s1) + s2**2)}, identical to the Core-15 Portfolio Std Dev PySR-only row: confirm which run wrote it.
(b) Seed-42 composition of the 23: the uploaded file gives 12 finite and 11 \textit{nan}; confirm no other copy was used. (The checker-asymmetry question this raised is closed by Item~18.)
(c) Hybrid catastrophic count is 0 under test $\Rsq<-10$; the attributed count is 5--7 per seed. (d) \texttt{generate\_figures.py} legacy
\texttt{SEED\_DATA} fallback (fix7) was not changed.
\end{itemize}


\subsection*{Item 17: Open points stripped from the manuscript (2026-10-08)}
\textbf{Status:} Open (tracking only). \textbf{Where:} \texttt{jmlr\_paper\_main.tex}, \texttt{supp\_benchmark\_report.tex},
\texttt{supp\_routing\_improvements.tex}. At the request of the author, every statement in the manuscript and supplements that
flagged an unresolved point, a pending rerun or an open question was removed or rewritten as a plain statement of what the paper reports.
The text removed, and what would close each point, is recorded here. The manuscript makes no claim that depends on any of these points.
\begin{enumerate}\small
\item \textbf{Core-15 tier criterion (main paper, ``Core-15 difficulty tiers'').} Removed: ``No formal criterion defines this assignment;
it is inherited from the original figure grouping in the figure-generation script.'' Now reads: a fixed row grouping used only to order the heatmaps,
not derived from a quantitative rule. \emph{To close:} state the rule that produced \texttt{\_HM\_TIER}, or confirm the grouping is purely presentational.
\item \textbf{Pointer to deferred items (main paper, note after the abstract).} Removed the clause ``the deferred items and their provenance are tracked in
the Post-Submission Report''. Also removed the unused \texttt{\textbackslash TODO} macro.
\item \textbf{Multi-run instability (main paper: Introduction footnote and paragraph, Michaelis--Menten case study, Hard-tasks paragraph,
\S``Stability Under Stochastic Inference'', Finding~1, Mode~A paragraph, design-principle bullet).} Removed: the statements that the $K=30$ sweep was planned but its
raw results are absent, that the $\II=0.0$ artifact exists and a rerun has not been performed, that the measurement is ``designed but not currently verified'', that the
regime table is withheld ``pending a genuine multi-run rerun'', that attribution to specific tasks ``is deferred to future work'', and that the Instability Index ``remains open''.
The paper now says only that it reports single-run results and no multi-run index. The underlying gap is Item~1 above and is unchanged.
\item \textbf{FIX-N3 action item (ii) (main paper, LLM-sampling paragraph).} Removed: ``A future re-benchmarked run should still record \texttt{temperature} and
\texttt{n\_candidates}\ldots; this remains open pending that rerun.'' The paper keeps the factual statement that \texttt{run\_all.sh} STEP~7/8 pass \texttt{--temperature 0.25}.
\emph{To close:} a full-budget rerun whose result JSON records both fields.
\item \textbf{M3 noiseless recovery (benchmark supplement, limitations list).} Removed: ``this should be treated as an open finding requiring investigation, not merely a timing/framing issue''
(M3 strict recovery is 23/30 at $\sigma=0$). The figure itself stays. \emph{To close:} investigate the seven M3 failures at $\sigma=0$.
\item \textbf{Process/CI credibility row and verdict (benchmark supplement, readiness table).} Removed the row marked OPEN
(``Baseline-lock now live-confirmed (74/90); do not cite July~9 dashboard until re-run'') and the verdict clause ``one process item above remains open''.
The dashboard remains uncitable (Item~C5); the supplement's limitations list still carries the resolved entry.
\item \textbf{Item~10b mechanism mismatch (benchmark supplement, timeout-upgrade note).} Removed the ``Update (Phase~3 audit, Item~10b, open)'' paragraph, the
description of the \texttt{\_ProcBox}/\texttt{\_kill\_process\_group} fix and its two wall-clock verifications, and ``see the open question above''. The mismatch:
the root cause is described as in-process PySR via \texttt{juliacall} (no separate OS process), while the fix threads a live \texttt{Popen} handle, which presupposes a
subprocess. Either the root-cause description is incomplete, the fix targets a different subprocess-based hang, or \texttt{verify\_proc\_box\_fix.py} exercises another path.
\emph{To close:} check the \texttt{\_ProcBox} implementation and the verification script's PySR invocation against \texttt{run\_comparative\_suite\_benchmark\_v2\_FIXED.py}.
The supplement keeps only the recorded 27574\,s wall-clock and the in-process explanation of the timeout gap.
\item \textbf{``Remaining Open Issues'' (routing supplement, former Sec.~Remaining Open Issues) and its pointer in the front note.} Removed the section. Its entries:
(1) formula-evaluator divergence --- closed by Fix~5; (2) stochastic NN initialisation for cases routed to \texttt{nn} (use a deterministic seed or model caching);
(3) probe threshold $\Delta R^2\ge0.15$ chosen heuristically (sweep $\{0.05,0.10,0.15,0.20,0.25\}$ on a held-out set); (4) probe sorts by the primary feature only
(a Mahalanobis-distance probe would be more robust); (5) two scores on cases 6--7 remain NaN from LLM API connection errors (retry with exponential backoff, $\le3$ attempts, delays 5/15/45\,s).
\item \textbf{Items not changed.} The ``Limitations and Future Work'' section of the main paper, the ``we have not rerun the baseline with the bypass disabled'' disclosures
(main paper and benchmark supplement), and the untested-domain notes in the supplement's limitations list are ordinary scope statements and were kept.
\end{enumerate}

\subsection*{Item 18: Paired checker audit and ``uncorroborated'' rewording (2026-10-08)}
\textbf{Question.} The pure-LLM checker (\texttt{test\_formula\_accuracy}) builds its namespace from numpy/math names and never reads
\texttt{metadata["constants"]}; the hybrid evaluator injects them. A formula using a bare constant name would therefore return
\texttt{NaN} in the pure-LLM arm but pass in hybrid, which could explain some of the 11 \texttt{NaN} cases among the 23.
\textbf{Test.} \texttt{paired\_checker\_audit.py} re-ran the pure-LLM arm on all 74 tasks (one LLM call per task) and scored each formula
three ways on the same split: A (paper's pure-LLM checker), B0 (hybrid evaluator, no constants), B1 (hybrid evaluator, with constants).
\textbf{Result.} 11 formulas failed under A; all 11 also failed under B0 and B1 (\texttt{formula\_fails\_everywhere}). Explained by missing
constants: 0; other evaluator difference: 0. Tasks above $\Rsq>0.99$: 48/74 under A, 47/74 under B1. An earlier 10-task trial (first 10 tasks, none with constants) gave 0 failures and 8/10 under both.
\textbf{Limits.} Formulas are re-sampled, so this measures how often the artifact occurs and does not identify the original 11 tasks;
the original formulas were not stored. It tests only the pure-LLM side: hybrid makes its own LLM call, so which arm's formula was better on any task is not determined.
The re-run's 48/74 versus the paper's 44/74 is sampling variance between two runs; the sample count/seed of the re-run is not recorded.
One task passes under A but fails under B1 (48 vs 47); this row of the full \texttt{paired\_audit.csv} has not been inspected, and until it is,
B1 should not be treated as a reference scorer.
\textbf{Manuscript changes.} (i) ``Fabricated'' replaced by ``uncorroborated'' throughout (definition added in the methodology paragraph of the main paper's hybrid decision-attribution section);
(ii) ``Composition of the 23'' rewritten to report the audit and to drop the sentence that the NaN cases may reflect the pure-LLM evaluator, and the $55/74=74.3\,\%$ sensitivity figure;
(iii) a limitation paragraph on independent sampling and checker asymmetry added to the Limitations section.
\textbf{Still open.} The ablation table (\texttt{tab:llm\_ablation}) is described as taken from \texttt{\_merged.json}, which is not in the repository, and it disagrees with the raw shards on at least four cells; the table has not been regenerated or disclosed as a different run.
Also unchanged: \texttt{tab:baseline} ($n{=}66$ denominator, no source file), reproducibility-checklist values, and the 25 tables marked UNVERIFIED.

\section{Closed Items (Confirmed, No Manuscript Change)}

\subsection*{Item C1: Hybrid decision-attribution bug uncorroborated-count (23)}
Independently recomputed from \texttt{hypatiax\_defi\_benchmark\_v3\_results\_seed42.json}
(canonical copy at that date, md5 \texttt{89037935469176741a5fb4479f2c725f}; now superseded by \texttt{cbf503d0\ldots}, same 23/3/11/9 result) using the
correct definition — the routed sub-method's own reported value must itself
support near-perfect, not merely carry a \texttt{success=True} flag. Result:
23 uncorroborated cases (3 easy / 11 medium / 9 hard), matching the manuscript's
the hybrid decision-attribution table task list exactly, name for name.
\textbf{No change made.} (An earlier, narrower recomputation using only the
boolean success flag gave 22 and was corrected after finding the
"Spot price from AMM reserve" case, where \texttt{hybrid} reports
$\Rsq{=}1.0$ against an actual routed \texttt{pure\_llm} result of
$\Rsq{=}0.981$ — a value-mismatch case the boolean-flag test missed.)
\textbf{Re-verified independently a second time (2026-09-27)} directly
against the same file retrieved fresh from
\texttt{github.com/sednabcn/LLM-HypatiaX-DEV}: 44/74 = 59.5\,\% corrected,
23 uncorroborated (Easy 3/24, Medium 11/29, Hard 9/21) — exact match.

\subsection*{Item C2: v3-vs-v4 methodology}
Confirmed: \texttt{hypatiax/data}'s \texttt{v4} result file uses a genuinely
different 4-arm router (\texttt{v4\_llm}, \texttt{v4\_linear\_fallback},
\texttt{v4\_residual\_nn}, \texttt{v4\_linear\_fallback\_local}), not a variant
of v3. The paper's v3 scope is legitimate and unaffected; the existing
one-line disclosure of active v4 development (main paper, end of
the hybrid decision-attribution section) is accurate as written. No change made.

\subsection*{Item C3: Redaction resolution pass across the two supplements
(2026-09-27)}
7 of 11 \texttt{[VALUE REDACTED]} placeholders (6 in the routing
supplement, 1 in the benchmark report) resolved by direct textual
cross-reference to this paper's own \S~``The Hybrid
Decision-Attribution Bug,'' with the source section cited inline at each
resolved location. No new data was introduced; every resolved value
(90.5\,\%, 98.6\,\%, 59.5\,\%) was already disclosed in this paper before
the pass. The remaining 4 placeholders and 1 newly-found figure conflict
are tracked as Open Item~2 above; they were not closed. See also Open
Item~3 above regarding the reliability of \texttt{paper\_audit.pdf}'s own
figures, found during the same pass.

\subsection*{Item C4: v4 attribution-flag bug — reverted to v3 (2026-09-27)}
While validating the v4 router (see Item C2), a bug was found in v4's own
decision-attribution flag: it can report \texttt{success=True} for a routed
sub-method even when that sub-method's own value does not actually meet the
near-perfect threshold — the same class of masking error identified and
fixed for v3 in Item C1, but reintroduced in v4's separate implementation.
Fixing it properly (auditing v4's routing logic the same way v3's was
audited) is out of scope for this submission. \textbf{Decision: revert to
v3 for all reported results and close this point for now.} v3's
attribution logic is the one already independently verified in Item C1; v4
remains under active development (as already disclosed in the manuscript)
and its attribution-flag bug should be fixed and re-audited before any v4
numbers are cited in a future revision.

\subsection*{Item C5: \texttt{build\_report.py} stale-green-dashboard fix — confirmed live (2026-09-27)}
The July~9 CI dashboard asserted ``0 open issues'' while its own upstream
notebooks (NB-01--NB-06) reported ``not found'' for that run --- flagged
as an open process item in both supplements and the main paper's
\S~``Outstanding Process/Credibility Items.'' Located
\texttt{scripts/build\_report.py} in the project repository and confirmed
the fix already exists and is live: \texttt{notebook\_execution\_status()}
checks every code cell's \texttt{execution\_count} (not mere file
existence) and, at the call site, collects any incomplete or missing
notebooks and calls \texttt{sys.exit(1)} before publishing the dashboard,
except on an explicit fix-only run. The function and call site are both
tagged \texttt{FIX-DASHBOARD-STALE-GREEN} in their own source comments and
cite ``verification report, item 10.1'' directly --- i.e. this project's
own pipeline was already patched in response to substantially this same
finding. Confirmed via \texttt{git log}: committed \texttt{9d7a833f},
2026-07-29, predating this review and postdating the July~9 incident it
addresses. \textbf{No further code change needed.} All three manuscript
documents updated (2026-09-27) to mark this item Resolved rather than
Open; the July~9 dashboard itself remains uncitable as a specific
artifact, but the underlying defect cannot recur silently.

\subsection*{Item C6: \texttt{paper\_audit.pdf} counts and cleanup-review problems 2--5 (2026-10-03)}
Evaluated against the current sources, with no manuscript change needed. Confirmed correct in the manuscript: $74/90$
(not $81/90$); the CI figures $[1087,1456]\%$, $1231\%$ and $n=14/13$ appear nowhere; the July~9 dashboard is flagged in all
three files; the $23$-case hybrid-bug count. Wrong or stale in \texttt{paper\_audit.pdf} (see Item~3 and Item~12): $34$
redactions ($0$ now), blank \texttt{main\_results} row (populated), $8$ uncorroborated cases ($23$), $2.5$--$9.3\times$ runtime range
(withdrawn), $0$ wired tables ($2$ \texttt{\textbackslash input}). Cleanup-review problems 2 (scalability speedups removed),
3 (NN time now $0.27$\,s pooled, $0.18$--$0.36$\,s per seed), 4 (complexity $1.4\times10^{8}$; $1.7\%$; $93.8\%$; both timing
blocks sum to $100\%$) and 5 (orphaned sentences) are fixed in the sources.

\subsection*{Item C7: Portfolio Variance seed-sweep data located and checked (2026-10-04)}
\texttt{portfolio\_variance\_seed\_sweep.json} (Exp-1 ablation layout) and five
\texttt{portfolio\_variance\_defi\_seed\_sweep\_seed<N>.json} shard files (v3c3, DeFi protocol) are now available.
\emph{Checked:} every far-$R^2$ in \texttt{tab:portfolio\_seed\_sweep} (both columns), both means ($-31.953$, $-6.261$), the
4/5 win count and the 2/5 exact-recovery count match the ablation file; the \texttt{[Unverified]} note in the caption was
replaced by a verification note. \emph{Not checked:} the Bayesian figure $P\approx0.76$ in the same section has no source in
either file. \emph{New finding:} the shard files report far-$R^2=1.000$ on 5/5 seeds, all with routing decision
\texttt{llm}; this is a different experiment from the table and was added to the manuscript as a separate paragraph, not
merged. \emph{Generator caveat:} \texttt{generate\_tables.py} merges the shard files with PySR-only values from the ablation
file, so a regenerated \texttt{portfolio\_sweep.tex} (5/5 exact) will not match the manuscript table (2/5 exact) and
would be reported DIFF by the table audit. Decide which source defines Table~5 before wiring it with
\texttt{\textbackslash input}. \texttt{generate\_figures.py} now reads the sweep JSON for the three portfolio figures; the
previously embedded values had PySR seed~2024 at far-$R^2=-12.11$ instead of $-118.448$.


\subsection*{Item C8: Noise-sensitivity tables checked against the sweep and raw files (2026-10-04)}
The five \texttt{noise\_sweep\_20260812\_*\_nshards0N.json} files and the five \texttt{protocol\_core} runs
(\texttt{noiseless\_20260812\_170344}, \texttt{noisy\_20260812\_172958}, \texttt{\_173425}, \texttt{\_173523}, \texttt{\_175318}) are now
available. \emph{Checked:} the noise levels (0, 0.0005, 0.001, 0.005, 0.01, i.e.\ 0, 0.05, 0.1, 0.5, 1\%; shard
\texttt{nshards01}--\texttt{05} map to 0, 0.05, 0.1, 0.5, 1\%); the run thresholds ($0.999999$, $0.95$, $0.95$, $0.995$, $0.99$) against the sweep
\texttt{threshold\_used} ($0.999999$, $0.995$, $0.995$, $0.995$, $0.99$), which is why the per-run and the \texttt{tab:rr\_noise}
tables differ at 0.05\% and 0.1\%; and every cell of both noise-count tables (per-run threshold, and fixed $0.99$), recomputed
independently from the raw runs and from the sweep \texttt{per\_equation} data (the two sources agree). Two cells were wrong and are corrected: SEL at
$\sigma=0$ $27/30\to28/30$ (pooled $146/150$) and PureLLM at $\sigma=0$ $29/30\to30/30$ (pooled $148/150$, in both tables).
\emph{Not rechecked:} \texttt{tab:r2\_noise}, \texttt{tab:rr\_noise}, \texttt{tab:time\_noise}. \emph{Still open:} the noise-sweep
shards do not store $n$; see Item~12 on the runner patch.

\subsection*{Item C9: \texttt{benchmark\_results.json} provenance and the SEL row (2026-10-04)}
The six-method \texttt{benchmark\_results.json} (180 records, 6 methods $\times$ 30 tests) was checked. At $\Rsq\ge0.999999$ it gives
PureLLM 29/30, ImprovedNN 0/30, M3 23/30, M4 30/30, SEL 27/30 (mean $\Rsq$ 0.999868, the figure in the removed footnote), HybridDiscovery
21/30. Compared test by test with the 170344 run: M3, M4 and HybridDiscovery are identical; ImprovedNN differs on all 30 tests but
stays at 0/30; SEL differs on four tests (Arrhenius, Nernst, Bose--Einstein, Fermi--Dirac), of which only Arrhenius ($0.99992$ there, $0.9999999987$ in 170344)
crosses the threshold, giving 27 versus 28; PureLLM differs on Curie's law ($0.0$ there, $1.0$ in 170344), giving 29 versus 30.
It is therefore a different run from 170344, and the supplement uses 170344 for every $\sigma=0$ cell. The script
\texttt{run\_comparative\_suite\_benchmark\_v2.py} (\texttt{print\_summary()}) appends to this file and drops only entries with the same (test, method), so
a file holds whatever methods were run into that path. \emph{Manuscript change:} the SEL row of \texttt{tab:overall} is now filled from 170344
(28/30, mean $\Rsq$ 0.999711, 39.1\,s) and the old ``Not reported'' footnote removed; the PureLLM row stays at 29/30 (per-domain recount, 18/19 live),
which the six-method file also reproduces (29/30, Curie's law failing). Main-table PureLLM 29/30 versus noise-table 30/30 is explained in the
supplement text. \emph{Not checked:} the 18/19 live-row split and the 11 hardcoded flags (the per-domain files were not supplied).
\emph{Code note:} the comment above \texttt{\_SHARD\_TAG} in the script says outputs get an \texttt{\_nshardsNN} suffix, but the code builds
\texttt{\_shrd<ID>} from \texttt{HYPATIAX\_SHARD\_ID}; the \texttt{\_nshardsNN} names in the noise sweep must come from the orchestrator.
Fix the comment or the code.

\end{document}
